%==============================================================================
%  UG PROJECT REPORT  --  Chandigarh University format
%  Project : An Interpretable Multi-Modal Spam Detection Framework Using
%            Explainable AI and Machine Learning
%  Compile with: pdfLaTeX  (Overleaf: Menu -> Compiler -> pdfLaTeX)
%  Compile 2-3 times so the Table of Contents / List of Figures fill in.
%
%  IMAGES live in the figures/ sub-folder (on Overleaf, upload the folder as is).
%  The document compiles even if one is missing (a labelled grey box is shown):
%     figures/cu_logo.png              <- Chandigarh University logo (cover+certificate)
%     figures/dataset_overview.png     <- class balance & message length (EDA)
%     figures/sms_roc_curves.png       <- ROC curves, SMS models
%     figures/email_roc_curves.png     <- ROC curves, e-mail models
%     figures/sms_shap_waterfall.png   <- SHAP waterfall of the exhibited SMS
%     figures/sms_shap_beeswarm.png    <- global SHAP beeswarm, SMS
%     figures/email_shap_beeswarm.png  <- global SHAP beeswarm, e-mail
%     figures/app_detector_tab.png     <- screenshot of the Streamlit Detector tab (Chapter 4)
%     figures/{sms,email}_model_comparison.png, {sms,email}_confusion_matrices.png,
%     figures/{sms,email}_lime_case_study.png  <- produced by src/train.py
%     figures/{sms,email}_shap_importance.png  <- produced by scripts/make_shap_figures.py
%  The three SHAP figures are generated by scripts/make_shap_figures.py.
%==============================================================================
\documentclass[12pt,a4paper]{report}

%------------------------- page geometry (1 inch all round) -------------------
\usepackage[a4paper,top=1in,bottom=1in,left=1in,right=1in]{geometry}
% For a binding margin instead, use:
% \usepackage[a4paper,top=1in,bottom=1in,left=1.25in,right=1in]{geometry}

%------------------------- fonts: Times New Roman clone -----------------------
\usepackage[T1]{fontenc}
\usepackage[utf8]{inputenc}
\usepackage{mathptmx}          % Nimbus Roman = metrically identical to Times New Roman

%------------------------- spacing, layout, tools -----------------------------
\usepackage{setspace}
\usepackage{graphicx}
\usepackage{xcolor}
\usepackage{titlesec}
\usepackage{tocloft}
\usepackage[font={bf,footnotesize},labelsep=colon,justification=centering,skip=6pt]{caption}
\usepackage{enumitem}
\usepackage{array}
\usepackage{longtable}
\usepackage{float}
\usepackage{multirow}
\usepackage{amsmath,amssymb}
\usepackage{pgfplots}                 % charts drawn directly in LaTeX
\pgfplotsset{compat=1.17}
\usepackage{pgfgantt}                 % Gantt chart for the timeline
\usetikzlibrary{arrows.meta,positioning,calc,shapes.geometric}
\usepackage[hidelinks,bookmarksnumbered]{hyperref}   % keep hyperref last

\onehalfspacing                        % 1.5 line spacing for the general text
\setlength{\parindent}{0pt}            % Word-style: no first-line indent
\setlength{\parskip}{6pt}
\setlist{nosep,topsep=4pt,itemsep=3pt,leftmargin=1.2cm}
\renewcommand{\arraystretch}{1.3}
% ragged-right paragraph column for tables:  L{3cm}  instead of  p{3cm}
\newcolumntype{L}[1]{>{\raggedright\arraybackslash}p{#1}}

% Word does not hyphenate by default -- switch it off so the justification matches
\frenchspacing          % Word-style: no extra space after a full stop
\hyphenpenalty=10000
\exhyphenpenalty=10000
\tolerance=9999
\emergencystretch=3em

%------------------------- exact font sizes from the format sheet -------------
\newcommand{\sixteen}{\fontsize{16}{19.2}\selectfont}   % chapter names
\newcommand{\fourteen}{\fontsize{14}{21}\selectfont}    % main headings
\newcommand{\twelve}{\fontsize{12}{18}\selectfont}      % body / sub-headings
\newcommand{\eighteen}{\fontsize{18}{27}\selectfont}    % report title

%------------------------- CHAPTER: 16pt, bold, centered ----------------------
\titleformat{\chapter}[display]
  {\filcenter\bfseries\sixteen}
  {CHAPTER \thechapter.}
  {6pt}
  {}
\titlespacing*{\chapter}{0pt}{-24pt}{22pt}

% unnumbered chapters (REFERENCES, APPENDIX, USER MANUAL ...)
\titleformat{name=\chapter,numberless}[display]
  {\filcenter\bfseries\sixteen}{}{0pt}{}
\titlespacing*{name=\chapter,numberless}{0pt}{-24pt}{22pt}

%------------------------- HEADING: 14pt, bold, left ---------------------------
\titleformat{\section}[hang]
  {\bfseries\fourteen}
  {\makebox[3.4em][l]{\thesection.}}{0pt}{}
\titlespacing*{\section}{0pt}{16pt}{8pt}

%------------------------- SUB-HEADING: 12pt, bold, left -----------------------
\titleformat{\subsection}[hang]
  {\bfseries\twelve}
  {\makebox[3.4em][l]{\thesubsection.}}{0pt}{}
\titlespacing*{\subsection}{0pt}{12pt}{6pt}

\titleformat{\subsubsection}[hang]
  {\bfseries\twelve}
  {\makebox[3.9em][l]{\thesubsubsection.}}{0pt}{}
\titlespacing*{\subsubsection}{0pt}{10pt}{6pt}

\setcounter{secnumdepth}{3}
\setcounter{tocdepth}{2}

%------------------------- TABLE OF CONTENTS styling ---------------------------
\renewcommand{\contentsname}{TABLE OF CONTENTS}
\renewcommand{\cfttoctitlefont}{\bfseries\sixteen}
\setlength{\cftbeforetoctitleskip}{0pt}
\setlength{\cftaftertoctitleskip}{22pt}

% "Chapter 1. ....... 4"
\renewcommand{\cftchappresnum}{Chapter\ }
\renewcommand{\cftchapaftersnum}{.}
\setlength{\cftchapnumwidth}{6.2em}
\renewcommand{\cftchapfont}{\bfseries}
\renewcommand{\cftchappagefont}{\bfseries}
\renewcommand{\cftchapleader}{\bfseries\cftdotfill{\cftdotsep}}
\renewcommand{\cftchapdotsep}{\cftdotsep}
\setlength{\cftbeforechapskip}{8pt}

\renewcommand{\cftsecfont}{\bfseries}
\renewcommand{\cftsecpagefont}{\bfseries}
\renewcommand{\cftsecleader}{\bfseries\cftdotfill{\cftdotsep}}
\setlength{\cftsecnumwidth}{3.2em}
\setlength{\cftbeforesecskip}{6pt}

\renewcommand{\cftsubsecfont}{\bfseries}
\renewcommand{\cftsubsecleader}{\bfseries\cftdotfill{\cftdotsep}}
\renewcommand{\cftsubsecpagefont}{\bfseries}
\setlength{\cftsubsecnumwidth}{4.0em}
\setlength{\cftbeforesubsecskip}{6pt}

%------------------------- LIST OF FIGURES / TABLES styling --------------------
\renewcommand{\listfigurename}{List of Figures}
\renewcommand{\cftloftitlefont}{\bfseries\sixteen}
\setlength{\cftbeforeloftitleskip}{0pt}
\setlength{\cftafterloftitleskip}{22pt}
\renewcommand{\cftfigpresnum}{Figure\ }
\renewcommand{\cftfigaftersnum}{:}
\setlength{\cftfignumwidth}{5.6em}
\renewcommand{\cftfigfont}{\bfseries}
\renewcommand{\cftfigpagefont}{\bfseries}
\renewcommand{\cftfigleader}{\bfseries\cftdotfill{\cftdotsep}}
\setlength{\cftbeforefigskip}{8pt}

\renewcommand{\listtablename}{List of Tables}
\renewcommand{\cftlottitlefont}{\bfseries\sixteen}
\setlength{\cftbeforelottitleskip}{0pt}
\setlength{\cftafterlottitleskip}{22pt}
\renewcommand{\cfttabpresnum}{Table\ }
\renewcommand{\cfttabaftersnum}{:}
\setlength{\cfttabnumwidth}{5.2em}
\renewcommand{\cfttabfont}{\bfseries}
\renewcommand{\cfttabpagefont}{\bfseries}
\renewcommand{\cfttableader}{\bfseries\cftdotfill{\cftdotsep}}
\setlength{\cftbeforetabskip}{8pt}

% centre the three list titles (tocloft's own \hfill trick loses to \parfillskip)
\makeatletter
\renewcommand{\@cftmaketoctitle}{%
  \addpenalty\@secpenalty \vspace*{\cftbeforetoctitleskip}\@cftpagestyle
  {\interlinepenalty\@M {\centering\cfttoctitlefont\contentsname\par}%
   \cftmarktoc \nobreak \vskip\cftaftertoctitleskip \@afterheading}}
\renewcommand{\@cftmakeloftitle}{%
  \addpenalty\@secpenalty \vspace*{\cftbeforeloftitleskip}\@cftpagestyle
  {\interlinepenalty\@M {\centering\cftloftitlefont\listfigurename\par}%
   \cftmarklof \nobreak \vskip\cftafterloftitleskip \@afterheading}}
\renewcommand{\@cftmakelottitle}{%
  \addpenalty\@secpenalty \vspace*{\cftbeforelottitleskip}\@cftpagestyle
  {\interlinepenalty\@M {\centering\cftlottitlefont\listtablename\par}%
   \cftmarklot \nobreak \vskip\cftafterlottitleskip \@afterheading}}
\makeatother

%------------------------- helper: logo with safe fallback ---------------------
\newcommand{\culogo}[1][2.6in]{%
  \IfFileExists{figures/cu_logo.png}
    {\includegraphics[width=#1]{figures/cu_logo.png}}
    {\fbox{\parbox[c][1.0in][c]{#1}{\centering\small\itshape
      upload \texttt{figures/cu\_logo.png}\\to the project folder}}}%
}

%------------------------- helper: image with safe fallback --------------------
% \figimg[width]{filename}{fallback-box-height}
\newcommand{\figimg}[3][0.75\textwidth]{%
  \IfFileExists{#2}
    {\includegraphics[width=#1]{#2}}
    {\fbox{\parbox[c][#3][c]{#1}{\centering\small\itshape
      upload \texttt{\detokenize{#2}}\\to the project folder}}}%
}

%------------------------- TikZ styles for the diagrams -------------------------
\tikzset{
  blk/.style={draw=black!70, fill=gray!8, rounded corners=2pt, align=center,
              font=\small, minimum height=0.85cm, inner sep=5pt},
  mdl/.style={blk, fill=blue!8},
  xai/.style={blk, fill=orange!14},
  dat/.style={blk, fill=green!10},
  arr/.style={-{Latex[length=2.6mm]}, semithick},
}

%------------------------- front-matter page helper ----------------------------
\newcommand{\frontchapter}[1]{%
  \chapter*{#1}\phantomsection\addcontentsline{toc}{chapter}{#1}%
}

%==============================================================================
%  >>>>>>>>>>>>>>>>>>>>>  EDIT YOUR DETAILS HERE  <<<<<<<<<<<<<<<<<<<<<<<<<<<<
%==============================================================================
\newcommand{\ProjectTitle}{AN INTERPRETABLE MULTI-MODAL SPAM DETECTION FRAMEWORK USING EXPLAINABLE AI AND MACHINE LEARNING}
\newcommand{\CandidateA}{Aariz (23BCS10422)}
\newcommand{\CandidateB}{Naman (23BCS11036)}
\newcommand{\CandidateC}{Vanshit Kumar (23BCS10087)}
\newcommand{\CandidateD}{Aditi Arya (23BCS11535)}
\newcommand{\CandidateE}{Avish (23BCS13898)}
\newcommand{\AllCandidates}{Aariz (23BCS10422), Naman (23BCS11036), Vanshit Kumar (23BCS10087), Aditi Arya (23BCS11535) and Avish (23BCS13898)}
\newcommand{\DegreeName}{BACHELOR OF ENGINEERING}
\newcommand{\BranchName}{COMPUTER SCIENCE AND ENGINEERING}
\newcommand{\MonthYear}{NOVEMBER 2026}                      % <-- EDIT month of submission
\newcommand{\HoDName}{Dr.~Surinder Kaur}                % <-- EDIT: fill in your HoD's name
\newcommand{\SupervisorName}{Sonam Juneja (E11002)}
\newcommand{\SupervisorDesignation}{Assistant Professor}    % <-- verify designation
\newcommand{\DeptName}{Department of Computer Science \& Engineering}
%==============================================================================

\begin{document}

%==============================================================================
%  COVER / TITLE PAGE
%==============================================================================
\pagestyle{empty}
\pagenumbering{roman}

\begingroup
\setlength{\parskip}{0pt}
\singlespacing
\begin{center}
\vspace*{0.12in}

{\bfseries\eighteen \ProjectTitle\par}

\vspace{0.34in}
{\bfseries\fourteen A PROJECT REPORT\par}

\vspace{0.30in}
{\bfseries\itshape\fourteen Submitted by\par}

\vspace{0.26in}
{\bfseries\sixteen \CandidateA\par}
\vspace{0.06in}
{\bfseries\sixteen \CandidateB\par}
\vspace{0.06in}
{\bfseries\sixteen \CandidateC\par}
\vspace{0.06in}
{\bfseries\sixteen \CandidateD\par}
\vspace{0.06in}
{\bfseries\sixteen \CandidateE\par}

\vspace{0.30in}
{\bfseries\itshape\fourteen in partial fulfillment for the award of the degree of\par}

\vspace{0.28in}
{\bfseries\sixteen \DegreeName\par}

\vspace{0.16in}
{\bfseries\fourteen IN\par}

\vspace{0.16in}
{\fourteen \BranchName\par}

\vspace{0.28in}
\culogo

\vspace{0.10in}
{\bfseries\fourteen Chandigarh University\par}

\vspace{0.22in}
{\fourteen \MonthYear\par}

\end{center}
\endgroup
\clearpage

%==============================================================================
%  BONAFIDE CERTIFICATE
%==============================================================================
\begin{center}
\vspace*{0.1in}
\culogo

\vspace{0.35in}
{\bfseries\sixteen BONAFIDE CERTIFICATE\par}
\end{center}

\vspace{0.35in}

{\fourteen
Certified that this project report \textbf{``\ProjectTitle''} is the bonafide
work of \textbf{``\AllCandidates''} who carried out the project work under
my/our supervision.\par}

\vspace{0.45in}

\begin{singlespace}
\noindent
\begin{minipage}[t]{0.47\textwidth}
\raggedright
\fourteen
\vspace*{0.45in}
\textbf{SIGNATURE}\\[0.18in]
\HoDName\\[0.10in]
\textbf{HEAD OF THE DEPARTMENT}\\[0.10in]
\DeptName
\end{minipage}
\hfill
\begin{minipage}[t]{0.47\textwidth}
\raggedright
\fourteen
\vspace*{0.45in}
\textbf{SIGNATURE}\\[0.18in]
\SupervisorName\\[0.10in]
\textbf{SUPERVISOR}\\[0.10in]
\SupervisorDesignation\\[0.10in]
\DeptName
\end{minipage}
\end{singlespace}

\vspace{0.45in}

{\fourteen Submitted for the project viva-voce examination held on
\rule{1.6in}{0.4pt}\par}

\vspace{0.45in}

\begin{singlespace}
\noindent
\begin{minipage}[t]{0.47\textwidth}
\raggedright\bfseries\fourteen INTERNAL EXAMINER
\end{minipage}
\hfill
\begin{minipage}[t]{0.47\textwidth}
\raggedright\bfseries\fourteen EXTERNAL EXAMINER
\end{minipage}
\end{singlespace}

\clearpage

%==============================================================================
%  ACKNOWLEDGEMENT
%==============================================================================
\begin{center}
\vspace*{0.1in}
{\bfseries\sixteen ACKNOWLEDGEMENT\par}
\end{center}
\vspace{0.3in}

We would like to express our sincere gratitude to our supervisor,
\SupervisorName, \SupervisorDesignation, \DeptName, Chandigarh University, for
the continuous guidance, constructive criticism and encouragement extended
throughout the course of this project. Her insistence that every design
decision be justified --- and every prediction be explainable --- shaped the
central theme of this work.

We are equally thankful to \HoDName, Head of the Department, for providing the
computing facilities and the academic environment that made this work
possible. Our thanks also go to the faculty members of the department for
their timely feedback during the project reviews.

Finally, we thank our families and friends for their constant support and
motivation.

\vspace{0.6in}
\hfill\begin{minipage}{0.5\textwidth}
\raggedright
\CandidateA\\
\CandidateB\\
\CandidateC\\
\CandidateD\\
\CandidateE
\end{minipage}

\clearpage

%==============================================================================
%  TABLE OF CONTENTS
%==============================================================================
\tableofcontents
\addtocontents{toc}{\protect\thispagestyle{empty}}
\clearpage

%==============================================================================
%  From here the roman page numbers become visible: i, ii, iii ...
%==============================================================================
\pagestyle{plain}
\setcounter{page}{1}

\phantomsection\addcontentsline{toc}{chapter}{List of Figures}
\listoffigures
\clearpage

\phantomsection\addcontentsline{toc}{chapter}{List of Tables}
\listoftables
\clearpage

%==============================================================================
%  ABSTRACT
%==============================================================================
\frontchapter{ABSTRACT}

Fraudulent messages now reach users through more than one channel, and the
machine-learning filters that intercept them are usually black boxes: they
return a verdict but no reason. Aggregate accuracy alone can conceal
systematic errors and dataset artefacts, leaves the end user with no basis to
contest a decision, and gives an analyst nothing to audit.

This project presents an interpretable multi-modal spam detection framework
covering both SMS and e-mail. Each modality is served by its own specialist
pipeline: text normalisation that maps high-signal entities to readable
placeholder tokens (\texttt{httpaddr}, \texttt{numbr}, \texttt{moneysymb},
\texttt{phonenumbr}), TF--IDF features over uni- and bi-grams combined with
fifteen engineered stylometric statistics, and a hybrid ensemble of Logistic
Regression, Random Forest and XGBoost fused by performance-weighted soft
voting. A lightweight router selects the appropriate specialist at inference
time. A dual explanation layer justifies every prediction: LIME explains the
complete ensemble locally, while exact TreeSHAP provides local and global
attributions for the boosted-tree component. The framework ships as a fully
reproducible Python pipeline and a four-tab interactive Streamlit application.

Evaluated on 17,158 messages from two public research corpora using an 80/20
stratified split and 5-fold cross-validation, the hybrid ensemble achieved
98.55\% accuracy (ROC-AUC 0.9941) on the UCI SMS Spam Collection and 98.25\%
accuracy (ROC-AUC 0.9983) on the Enron-Spam corpus, matching (to four
decimal places) or beating every base learner on ROC-AUC. Beyond the numbers, the explanation layer exposed a
genuine dataset bias --- the token \texttt{enron} acting as a shortcut for
legitimate mail --- that no accuracy table could reveal, demonstrating the
practical value of shipping explanations together with predictions.

\vspace{12pt}
\textbf{Keywords:} spam detection, explainable AI, ensemble learning, SHAP,
LIME, TF--IDF, multi-modal classification, machine learning.

\clearpage

%==============================================================================
%  GRAPHICAL ABSTRACT
%==============================================================================
\frontchapter{GRAPHICAL ABSTRACT}

\vspace{0.5in}
\begin{center}
\begin{tikzpicture}[node distance=0.55cm and 0.85cm]
  \node[dat, minimum width=2.6cm] (msg) {SMS / E-mail\\message};
  \node[blk, right=of msg, minimum width=2.8cm] (feat)
       {TF--IDF (1--2 grams)\\+ 15 stylometric\\features};
  \node[mdl, right=of feat, minimum width=3.1cm] (ens)
       {Hybrid ensemble\\LR + RF + XGBoost\\(weighted soft vote)};
  \node[blk, right=of ens, minimum width=2.6cm, fill=green!12] (out)
       {Verdict +\\spam probability};
  \node[xai, below=of ens, minimum width=3.1cm] (xai)
       {LIME + SHAP\\``\emph{why} was this\\flagged?''};
  \draw[arr] (msg) -- (feat);
  \draw[arr] (feat) -- (ens);
  \draw[arr] (ens) -- (out);
  \draw[arr] (ens) -- (xai);
  \draw[arr] (feat.south) |- (xai.west);
\end{tikzpicture}
\end{center}

\vspace{0.35in}
Every message is routed to a modality-specific feature space, scored by a
hybrid ensemble of three heterogeneous learners, and returned to the user
together with a human-readable explanation of the verdict.

\clearpage

%==============================================================================
%  MAIN TEXT
%==============================================================================
\pagenumbering{arabic}
\setcounter{page}{1}

%------------------------------------------------------------------------------
\chapter{INTRODUCTION}

\section{Client Identification/Need Identification/Identification of relevant Contemporary issue}

Unsolicited and fraudulent messaging is a documented, growing contemporary
problem. The FBI's Internet Crime Complaint Center recorded close to
3,00,000 phishing/spoofing complaints in calendar year 2023 --- the single
largest crime category in its annual report --- and industry telemetry has for
years placed unsolicited traffic at roughly half of all e-mail exchanged
worldwide. The problem is equally visible on the SMS channel: as banking,
commerce and one-time-password flows moved to text messaging, so did the
scams (``smishing''). In India, the volume of unsolicited commercial
communication led the Telecom Regulatory Authority of India to notify the
Telecom Commercial Communications Customer Preference Regulations, 2018, and
to mandate a distributed-ledger registration system for all bulk senders ---
regulatory evidence that the issue is both real and unresolved.

The most recent edition of the same report, covering 2025, shows the problem
still growing: the Internet Crime Complaint Center received more than one
million complaints with reported losses of US\$20.9 billion, phishing and
spoofing was again the most reported category (191,561 complaints), and
22,364 complaints involved artificial intelligence (Federal Bureau of
Investigation, 2026). Generative AI is changing the character of the threat as
well as its volume: large language models can now write persuasive,
personalised phishing at negligible cost (Heiding et al., 2024; Yao et al.,
2024), and phishing produced in this way has been shown to get past the filters
of major webmail providers (Opara et al., 2025).

Filters exist, but a second, newer need has emerged around \emph{transparency}.
Modern spam classifiers are statistical black boxes: a message disappears and
nobody --- not the end user, not the platform analyst, not the auditor --- can
say why. Transparency obligations in data-protection law (Article~22 of the
GDPR in Europe; the Digital Personal Data Protection Act, 2023 in India) are
pushing automated decision systems toward explainability, and a filter that
cannot justify its verdicts is increasingly difficult to defend.

The client for this work is any operator of a user-facing messaging channel
--- typified by a university IT cell that receives phishing-report tickets
from students and staff --- who needs a filter that is (a)~accurate across
more than one message channel, (b)~able to attach a human-readable
justification to every verdict so that flagged messages can be triaged and
contested, and (c)~cheap enough to train and audit on commodity hardware.
This is a consultancy-style need: the client does not merely want a spam
score, they want a \emph{defensible} spam score.

\section{Identification of Problem}

No openly available framework simultaneously (i)~handles both SMS and e-mail,
whose statistics differ radically (median message length of 12 words versus
155; spam prevalence of about 12\% versus about 48\% in the public
research corpora), (ii)~matches the state-of-the-art accuracy of classical
machine-learning filters, and (iii)~attaches a faithful, per-prediction
explanation to every verdict. The consequences of this gap are concrete:
verdicts cannot be verified by non-experts, misclassified legitimate messages
cannot be contested, and hidden dataset biases pass silently into deployed
models. The problem addressed by this project is therefore the design,
implementation and validation of an accurate, fully interpretable,
multi-modal spam detection framework.

\section{Identification of Tasks}

The work was decomposed into the following tasks:

\begin{enumerate}
  \item Survey of the literature on spam/scam detection and on post-hoc
        explainable-AI techniques.
  \item Acquisition, cleaning and de-duplication of two public corpora (UCI
        SMS Spam Collection; Enron-Spam), followed by exploratory data
        analysis.
  \item Design of an interpretability-preserving preprocessing scheme and
        feature space (readable placeholder tokens; engineered stylometric
        features).
  \item Training and tuning of three heterogeneous base learners per modality
        with class-imbalance handling.
  \item Design and implementation of the performance-weighted soft-voting
        ensemble and the modality router.
  \item Integration of the explanation layer (LIME for the full ensemble;
        exact TreeSHAP for the boosted-tree component) at local and global
        scope.
  \item Development of the four-tab interactive Streamlit demonstration
        application and an automated end-to-end smoke-test suite.
  \item Evaluation on held-out test sets, error and bias analysis, and
        documentation of the results in this report.
\end{enumerate}

\section{Timeline}

The project was executed over one semester of approximately fifteen working
weeks. The schedule is represented as a Gantt chart in
Figure~\ref{fig:gantt}; the model-training, explanation-layer and application
workstreams deliberately overlap between weeks 8 and 12 because the
explanation artefacts were used to debug the models as they were being
trained.

\begin{figure}[H]
  \centering
  \resizebox{0.98\textwidth}{!}{%
  \begin{ganttchart}[
      x unit=0.62cm, y unit title=0.62cm, y unit chart=0.56cm,
      canvas/.append style={draw=black!60},
      title/.append style={draw=black!60, fill=gray!12},
      title label font=\scriptsize\bfseries,
      bar/.append style={fill=gray!35, draw=black!70},
      bar height=0.55,
      bar label font=\scriptsize
    ]{1}{15}
    \gantttitle{Project week}{15}\\
    \gantttitlelist{1,...,15}{1}\\
    \ganttbar{Literature survey}{1}{3}\\
    \ganttbar{Data acquisition \& EDA}{2}{4}\\
    \ganttbar{Preprocessing \& feature design}{4}{6}\\
    \ganttbar{Model training \& ensembling}{6}{9}\\
    \ganttbar{XAI layer (SHAP + LIME)}{8}{11}\\
    \ganttbar{Streamlit demo application}{10}{12}\\
    \ganttbar{Validation \& smoke testing}{12}{14}\\
    \ganttbar{Report writing}{13}{15}
  \end{ganttchart}}
  \caption{Gantt chart of the project schedule}
  \label{fig:gantt}
\end{figure}

\section{Organization of the Report}

Chapter~2 reviews the reported literature on spam detection and explainable
AI and closes with the formal problem definition and objectives. Chapter~3
covers the evaluation of candidate features, the design constraints, the two
alternative design flows that were considered and the justification for the
selected design. Chapter~4 presents the implementation using modern
engineering tools together with the complete experimental results, error
analysis and explainability findings. Chapter~5 concludes the work and
identifies the way ahead. The appendices list the software environment, the
complete engineered feature set, the derived ensemble weights and the
scam-keyword lexicon.

%------------------------------------------------------------------------------
\chapter{LITERATURE REVIEW/BACKGROUND STUDY}

\section{Timeline of the reported problem}

Unsolicited electronic messaging is almost as old as electronic messaging
itself: the first documented mass advertisement was sent over ARPANET in
1978. Volumes exploded with commercial e-mail in the 1990s, and the first
generation of defences was rule-based --- keyword blacklists, sender
blocklists and hand-written heuristics. The statistical era opened when
Sahami et al.\ (1998) framed junk-mail filtering as Naive-Bayes text
classification, followed quickly by margin-based methods (Drucker et al.,
1999) and, through the 2000s, by tree ensembles such as Random Forests
(Breiman, 2001). Research on the SMS channel became reproducible when
Almeida and G\'{o}mez Hidalgo (2011) released the UCI SMS Spam Collection,
while the Enron-Spam corpora (Metsis et al., 2006) did the same for
personalised e-mail filtering. From roughly 2015, deep neural models pushed
benchmark accuracy still higher but at a severe cost in transparency ---
which in turn triggered the explainable-AI wave: LIME (Ribeiro et al., 2016)
and SHAP (Lundberg and Lee, 2017) made post-hoc, per-prediction explanation
practical. Most recently, regulation (GDPR, 2018; DPDP Act, 2023) has begun
to treat explainability as a compliance requirement rather than a luxury,
which is the point in the timeline where this project positions itself.

\section{Proposed solutions}

Reported approaches fall into four broad families.

\subsection{Rule-based and collaborative filtering}

Blacklists, keyword rules and community-scored heuristic engines (e.g.\
SpamAssassin-style scoring). They are fully transparent and cheap, but
brittle: spammers evade fixed rules within days, false-positive rates are
high, and the rule bases require continuous manual curation.

\subsection{Classical statistical learning}

Naive Bayes, logistic regression, support vector machines and tree ensembles
trained on bag-of-words or TF--IDF representations. These methods reach
95--99\% accuracy on the public benchmark corpora, train in seconds to
minutes on commodity CPUs, and expose at least coarse interpretability
(coefficients, feature importances). Class imbalance --- severe on the SMS
channel --- must be handled explicitly.

\subsection{Deep learning}

Convolutional, recurrent and, more recently, transformer-based text
classifiers. They deliver the highest raw benchmark numbers but demand GPUs
and large training sets, and their decision process is opaque; attention
weights are known to be an unreliable proxy for faithful explanation.

\subsection{Explainable-AI methods}

Post-hoc techniques that explain an already-trained classifier. LIME fits a
local linear surrogate around one prediction of \emph{any} model by
perturbing the input; SHAP computes additive Shapley-value attributions, with
an exact polynomial-time algorithm (TreeSHAP) for tree ensembles. Combining a
strong classical ensemble with these two techniques promises near
state-of-the-art accuracy \emph{with} faithful explanations.

\subsection{Recent advances (2024--2026)}

Research on message filtering has moved quickly in the last three years. The
studies most relevant to this project are summarised below by theme.

\textbf{SMS spam and smishing.} Deep and transformer models dominate recent
SMS work. Mehmood et al.\ (2024) designed a deep-learning smishing detector
aimed explicitly at reducing false positives, and Uddin et al.\ (2025)
combined a transformer language model with an explainability analysis. Prince
et al.\ (2026) proposed an explainable hybrid transformer that fuses several
feature types for smishing detection, while Jasim et al.\ (2026) paired fuzzy
logic, which uses Gaussian membership functions to capture linguistic
uncertainty, with a bidirectional LSTM to make SMS filtering transparent. Hur
et al.\ (2026) benchmarked large language models (LLMs) on English and Korean
SMS datasets with zero-shot, few-shot and chain-of-thought prompting and found
them robust under concept drift and adversarial attack. Two further findings
are especially relevant here. \c{S}ahin et al.\ (2026) evaluated 42
combinations of seven text representations and six classical classifiers on
the same SMS corpus and concluded that the feature representation matters
more than classifier complexity, with character 3-grams and Logistic
Regression reaching 98.55\% accuracy. Salman et al.\ (2024) released a new SMS
dataset and analysed the evasive techniques that spammers use against
filters, and Mambina et al.\ (2024) extended deep SMS detection to Swahili.

\textbf{E-mail spam and phishing.} Fine-tuned transformers and LLMs are now
the most common e-mail detectors: Jamal et al.\ (2024) built an improved
transformer for three-way phishing/spam/ham classification, Roumeliotis et
al.\ (2024) compared fine-tuned LLMs, NLP models and CNNs, Uddin et al.\ (2026)
explained a transformer phishing detector, and Koide et al.\ (2026) used an
LLM to detect phishing and explain its verdicts to end users. The margin over
classical methods is small: Mel\'{e}ndez et al.\ (2024) found that RoBERTa
reached 0.9943 accuracy on phishing e-mail against 0.9876 for the best
traditional model, a support vector machine. Ensembles remain competitive:
Kshirsagar et al.\ (2025) reported 0.9905 accuracy and 0.9991 ROC-AUC with a
meta-learner on Enron-Spam combined with TREC 2007, and Hadi et al.\ (2026)
fused shallow lexical features with deep contextual embeddings in a hybrid
aimed at both channels, which they evaluated on the SMS Spam Collection
(98.8\% accuracy, ROC-AUC 0.993). On the attack side, Heiding et al.\ (2024)
showed that GPT-4 can generate phishing e-mails automatically, and Opara et
al.\ (2025) found that AI-generated phishing got past the filters of major
webmail providers but could be detected with 60 stylometric features and
XGBoost (96\% accuracy) --- direct support for the stylometric features used
in this project.

\textbf{Explainability in message filtering.} Shambour et al.\ (2025) combined
metaheuristic feature selection with SHAP explanations for spam detection,
and Panthi et al.\ (2026) compared SHAP and LIME on three boosting models for
phishing URLs, finding SHAP's attributions more stable at the global level and
LIME's more instance-specific. Several authors warn that explanations cannot
be taken at face value. Lyu et al.\ (2024) reviewed more than 110 explanation
methods in NLP through the lens of faithfulness, Krishna et al.\ (2024) showed
that different explanation methods frequently disagree and that practitioners
lack a principled way to resolve the disagreement, and Longo et al.\ (2024)
set out the open challenges of the field.

\textbf{Shortcut learning.} A model can reach high accuracy by exploiting
spurious cues that separate the classes in one dataset instead of learning
the intended concept. Du et al.\ (2024) documented such shortcuts in large
language models, and Dogra et al.\ (2024) surveyed how explanation methods can
expose the dataset biases behind them. Spam corpora are particularly exposed
because their legitimate messages often come from a single organisation ---
precisely the situation this project later uncovered in the Enron-Spam
corpus.

\section{Bibliometric analysis}

A comparison of the reviewed families on key features, effectiveness and
drawbacks is summarised in Table~\ref{tab:biblio}.

\begin{table}[H]
  \centering
  \caption{Comparison of reported spam-detection approaches}
  \label{tab:biblio}
  \begin{tabular}{|L{3.1cm}|L{3.4cm}|L{3.1cm}|L{4.0cm}|}
    \hline
    \textbf{Approach} & \textbf{Key feature} & \textbf{Effectiveness} & \textbf{Drawback} \\
    \hline
    Rule-based filters & Hand-written, fully transparent rules & Moderate; high false positives & Easily evaded; constant manual updating \\
    \hline
    Naive Bayes & Fast probabilistic word model & 95--98\% on benchmark corpora & Feature-independence assumption \\
    \hline
    SVM / Logistic Regression & Margin / linear decision boundary & 97--98\% & Linear view of language; needs calibration \\
    \hline
    Tree ensembles (RF, XGBoost) & Non-linear feature interactions & 97--99\% & Interpretability limited to importances \\
    \hline
    Deep / transformer models & Learned contextual representations & 98--99\%+ & Opaque; GPU and data hungry \\
    \hline
    Classical ensemble + XAI (this work) & Accuracy of ensembles with per-verdict explanations & 98.3--98.6\% (measured) & Explanation adds inference cost \\
    \hline
  \end{tabular}
\end{table}

Table~\ref{tab:recent} places the recent studies reviewed above next to this
project. Their figures come from different datasets and protocols, so they
indicate the state of the art rather than allowing a strict ranking.

\begin{table}[!htbp]
  \centering
  \caption{Recent studies (2024--2026) related to this project}
  \label{tab:recent}
  \small
  \begin{tabular}{|L{2.8cm}|L{2.2cm}|L{3.0cm}|L{2.4cm}|L{2.9cm}|}
    \hline
    \textbf{Study} & \textbf{Channel} & \textbf{Approach} & \textbf{Explanation} & \textbf{Reported result} \\
    \hline
    \c{S}ahin et al.\ (2026) & SMS & 42 pairs of classical representations and classifiers & -- & 98.55\% accuracy (character 3-grams + LR) \\
    \hline
    Hadi et al.\ (2026) & SMS (e-mail in scope) & Deep--shallow hybrid & -- & 98.8\% accuracy, ROC-AUC 0.993 \\
    \hline
    Hur et al.\ (2026) & SMS (English, Korean) & LLM prompting & LLM rationales & Robust under drift and attack \\
    \hline
    Mel\'{e}ndez et al.\ (2024) & E-mail (phishing) & Transformers vs.\ traditional ML & -- & RoBERTa 0.9943 vs.\ SVM 0.9876 accuracy \\
    \hline
    Kshirsagar et al.\ (2025) & E-mail & Meta-learner ensemble & Interpretable models & 0.9905 accuracy, 0.9991 ROC-AUC \\
    \hline
    Opara et al.\ (2025) & E-mail (phishing written by AI) & 60 stylometric features + XGBoost & -- & 96\% accuracy \\
    \hline
    This project & SMS + e-mail (routed) & TF--IDF + stylometry; weighted soft vote of LR, RF, XGBoost & LIME + exact TreeSHAP & 98.55\% / 98.25\% accuracy; ROC-AUC 0.9941 / 0.9983 \\
    \hline
  \end{tabular}
\end{table}

\section{Review Summary}

The literature exhibits a persistent trade-off: the most transparent systems
are the least accurate, and the most accurate are the least transparent.
Furthermore, nearly all published work targets a single channel, even though
SMS and e-mail scams are two faces of the same problem with radically
different message statistics. A hybrid ensemble of heterogeneous classical
learners, wrapped in a dual model-agnostic/model-specific explanation layer
and specialised per modality, occupies the unclaimed middle ground: benchmark
accuracy within a point of deep models, full per-prediction transparency, and
CPU-only training. This project adopts exactly that position.

The recent literature sharpens three further gaps. First, in the studies
reviewed here explanations are mostly presented as illustrations of single
predictions, and their faithfulness is rarely tested, although the XAI
literature shows that it cannot be assumed. Second, high benchmark accuracy is
seldom examined for shortcuts, even though organisation-specific corpora such
as Enron-Spam invite them. Third, model comparisons on near-saturated
benchmarks seldom report confidence intervals or significance tests, so
differences in the third decimal place are hard to interpret. This project
addresses all three in addition to the two-channel design.

\section{Problem Definition}

To design, build and validate an interpretable multi-modal spam detection
framework that classifies both SMS and e-mail messages as legitimate or
fraudulent using a hybrid ensemble of Logistic Regression, Random Forest and
XGBoost, and that attaches a faithful local explanation (LIME over the full
ensemble; SHAP over the boosted-tree component) to every verdict, together
with global model summaries. The framework must be reproducible end-to-end on
a commodity CPU and be operable by a non-expert through an interactive
application. The work explicitly excludes non-English corpora, adversarially
obfuscated spelling, image-based spam, and hardening for production
deployment.

\section{Goals/Objectives}\label{sec:objectives}

\begin{itemize}
  \item Achieve at least 97\% test-set accuracy on both public corpora (UCI
        SMS Spam Collection; Enron-Spam).
  \item Ensure the hybrid ensemble matches or beats every base learner on
        ROC-AUC for both modalities.
  \item Attach a local explanation to every single prediction and provide
        global explanations of what each model has learned.
  \item Keep the complete pipeline --- data download, training, evaluation,
        figure generation --- reproducible on a commodity CPU in under five
        minutes, verified by an automated test suite.
  \item Deliver an interactive demonstration application usable by a
        non-expert.
\end{itemize}

%------------------------------------------------------------------------------
\chapter{DESIGN FLOW/PROCESS}

\section{Evaluation \& Selection of Specifications/Features}

The candidate features identified in the literature were critically evaluated
against the client need of \emph{accurate and auditable} filtering.

\textbf{Retained.} Modality-specific pipelines (one specialist per channel,
selected by a lightweight, user-overridable router); normalisation of
high-signal entities into \emph{readable} placeholder tokens
(\texttt{httpaddr}, \texttt{numbr}, \texttt{moneysymb}, \texttt{phonenumbr})
instead of deleting them; fifteen engineered stylometric features computed on
the raw message and named in UPPER-CASE so that they remain visually distinct
from vocabulary in every explanation plot; TF--IDF over uni- and bi-grams
with sublinear term frequency; a heterogeneous three-model ensemble
(linear~/ bagging~/ boosting) fused by performance-weighted soft voting;
balanced class weights for imbalance handling; and a dual explanation layer
(LIME + SHAP).

\textbf{Rejected.} Stemming and lemmatisation (they destroy the surface word
forms on which readable explanations depend); deletion of URLs, numbers and
currency symbols during cleaning (they are among the strongest scam signals);
deep contextual models as the primary classifier (opaque, GPU-dependent);
synthetic minority oversampling of text (it manufactures training
``messages'' that never existed and cannot be meaningfully explained); and an
on-premise mail-server integration, which added engineering cost without
addressing the stated research need.

The retained normalisation step is specified in Table~\ref{tab:normalise}. The
rules are applied in the order listed, after the text has been lower-cased, so
that, for example, a phone number is recognised before its digits could be
replaced by the generic number token.

\begin{table}[!htbp]
  \centering
  \caption{Normalisation rules applied before TF--IDF vectorisation (in order)}
  \label{tab:normalise}
  \begin{tabular}{|c|L{5.8cm}|L{2.4cm}|L{4.4cm}|}
    \hline
    \textbf{Step} & \textbf{Pattern} & \textbf{Replaced by} & \textbf{Reason} \\
    \hline
    1 & HTML tags & a space & Markup carries no meaning for the reader \\
    \hline
    2 & URLs (\texttt{http(s)://\ldots}, \texttt{www.\ldots}, or a word ending in \texttt{.com}, \texttt{.net}, \texttt{.org}, \texttt{.info}, \texttt{.co.uk}) & \texttt{httpaddr} & Links are a primary scam signal \\
    \hline
    3 & E-mail addresses & \texttt{emailaddr} & Contact details in unsolicited messages \\
    \hline
    4 & A currency symbol (dollar, pound, euro, rupee) followed by an amount, or an amount followed by a currency word (dollars, pounds, euros, rs, \ldots) & \texttt{moneysymb} & Prizes, fees and refunds \\
    \hline
    5 & Phone numbers (8--15 digits with optional \texttt{+} and separators) & \texttt{phonenumbr} & Premium-rate and call-back numbers \\
    \hline
    6 & Remaining numbers & \texttt{numbr} & Keeps the signal without inflating the vocabulary \\
    \hline
    7 & Other non-letter characters & removed & Leaves clean, readable word tokens \\
    \hline
  \end{tabular}
\end{table}

\section{Design Constraints}

\begin{itemize}
  \item \textbf{Regulations:} transparency expectations of GDPR Article~22
        and the DPDP Act, 2023 motivate per-prediction explanations; the
        research licences of both corpora restrict the system to research and
        education use, and it is labelled accordingly.
  \item \textbf{Economic:} zero licence cost (fully open-source stack) and
        commodity-CPU training; no GPU may be assumed.
  \item \textbf{Environmental:} CPU-only training completing in about two
        minutes keeps the energy footprint negligible compared with
        GPU-trained deep models.
  \item \textbf{Health and Safety (user safety):} a false positive can
        suppress a legitimate OTP or an important mail, so high precision is
        prioritised on the SMS channel; no live user data is collected ---
        only public, anonymised research corpora are processed.
  \item \textbf{Manufacturability (reproducibility):} pinned dependencies,
        fixed random seeds, a single-command training pipeline and a
        13-check automated smoke-test suite guarantee that any examiner can
        rebuild the system.
  \item \textbf{Ethical and Social:} the scam-keyword lexicon is a plain,
        inspectable list in \texttt{src/config.py}; global SHAP summaries are
        used to audit the models for dataset bias; both corpora are credited
        to their original authors.
\end{itemize}

\section{Analysis and Feature finalization subject to constraints}

The constraint set directly shaped the final feature list. The
interpretability requirement eliminated stemming and deep contextual
encoders, and dictated the readable-placeholder normalisation: a token such
as \texttt{httpaddr} can appear verbatim inside a LIME highlight, whereas a
deleted URL cannot appear anywhere. The commodity-CPU constraint fixed the
representation at TF--IDF plus engineered statistics, giving compact feature
spaces of 3{,}015 dimensions for SMS and 5{,}015 for e-mail. The user-safety
constraint resolved the imbalance-handling question in favour of balanced
class weights and \texttt{scale\_pos\_weight} rather than synthetic
oversampling, so that every training example remains a real, explainable
message. Finally, the reproducibility constraint required every artefact ---
processed datasets, trained detectors, metrics, figures --- to be regenerated
by one command.

\section{Design Flow}

Two alternative designs were carried forward for detailed comparison.

\subsection{Design~I: Fine-tuned transformer classifier}

A single compact transformer (e.g.\ DistilBERT) fine-tuned jointly on both
corpora, with attention maps and gradient-based attributions offered as
explanations. Expected accuracy is at or marginally above the classical
state of the art, and one model serves both channels. However, fine-tuning
requires GPU time, inference is an order of magnitude slower on CPU, and ---
decisively --- attention weights are an unfaithful basis for explanation, so
the design fails the core transparency requirement of the client.

\subsection{Design~II: Hybrid classical ensemble with a dual XAI layer}

Each modality receives its own feature space and its own trained ensemble of
three heterogeneous learners --- Logistic Regression, Random Forest and
XGBoost --- combined in probability space with weights proportional to each
model's cross-validated F1. LIME explains the complete ensemble locally by
re-scoring word-dropout perturbations through the whole pipeline; exact
TreeSHAP explains the XGBoost component locally and globally. Training is
CPU-only and completes in about two minutes. The block diagram of this
alternative is shown in Figure~\ref{fig:block}.

\begin{figure}[H]
  \centering
  \begin{tikzpicture}[node distance=0.45cm and 0.7cm]
    \node[dat, minimum width=4.2cm] (raw)  {Raw message (SMS / E-mail)};
    \node[blk, below=of raw, minimum width=4.2cm] (rout) {Modality router\\(length/structure heuristic, user-overridable)};
    \node[blk, below=of rout, minimum width=7.6cm] (norm)
      {Normalisation: URLs $\to$ \texttt{httpaddr} $\cdot$ numbers $\to$ \texttt{numbr} $\cdot$ money $\to$ \texttt{moneysymb}};
    \node[blk, below=of norm, minimum width=7.6cm] (feat)
      {Feature extraction: TF--IDF (1--2 grams) + 15 engineered stylometric features};
    \node[mdl, below=1.0cm of feat] (rf) {Random\\Forest};
    \node[mdl, left=of rf]  (lr) {Logistic\\Regression};
    \node[mdl, right=of rf] (xgb) {XGBoost};
    \node[blk, below=1.0cm of rf, minimum width=6.4cm, fill=blue!12] (ens)
      {Hybrid soft-voting ensemble\\(weights $\propto$ cross-validated F1)};
    \node[dat, below left=0.75cm and -1.3cm of ens, minimum width=3.6cm, fill=green!12] (ver)
      {Verdict +\\spam probability};
    \node[xai, below right=0.75cm and -1.3cm of ens, minimum width=4.4cm] (xai)
      {XAI layer:\\LIME (full ensemble)\\SHAP (exact TreeSHAP)};
    \draw[arr] (raw) -- (rout);
    \draw[arr] (rout) -- (norm);
    \draw[arr] (norm) -- (feat);
    \draw[arr] (feat.south) -- (lr.north);
    \draw[arr] (feat.south) -- (rf.north);
    \draw[arr] (feat.south) -- (xgb.north);
    \draw[arr] (lr.south)  -- (ens.north);
    \draw[arr] (rf.south)  -- (ens.north);
    \draw[arr] (xgb.south) -- (ens.north);
    \draw[arr] (ens.south) -- (ver.north);
    \draw[arr] (ens.south) -- (xai.north);
    \draw[arr] (feat.east) .. controls +(2.2,0) and +(1.9,1.4) .. (xai.east);
  \end{tikzpicture}
  \caption{Block diagram of the proposed interpretable multi-modal detector (Design~II)}
  \label{fig:block}
\end{figure}

\section{Design selection}

The two alternatives were scored against the constraint set, as shown in
Table~\ref{tab:designsel}. Design~II was selected: it reaches the accuracy
band of the transformer on these corpora while satisfying every constraint
the transformer violates, and its tree-based component admits \emph{exact}
Shapley-value explanations rather than approximations.

\begin{table}[H]
  \centering
  \caption{Comparison and selection of the design alternatives}
  \label{tab:designsel}
  \begin{tabular}{|L{4.4cm}|L{4.3cm}|L{4.6cm}|}
    \hline
    \textbf{Criterion} & \textbf{Design~I (transformer)} & \textbf{Design~II (hybrid + XAI)} \\
    \hline
    Expected / measured accuracy & 98--99\% (reported) & 98.25--98.55\% (measured) \\
    \hline
    Explanation fidelity & Low (attention is unfaithful) & High (model-agnostic LIME; exact TreeSHAP) \\
    \hline
    Training hardware / time & GPU; tens of minutes--hours & CPU; $\approx$2 minutes \\
    \hline
    Reproducibility for examiners & Moderate & High (single command + tests) \\
    \hline
    Meets full constraint set & No & Yes \\
    \hline
  \end{tabular}
\end{table}

\section{Implementation plan/methodology}

The selected design is realised by the pipeline of Figure~\ref{fig:flow}.
Text is vectorised with sublinear TF--IDF using the smoothed inverse document
frequency of scikit-learn,
\begin{equation}
  w_{t,d} \;=\; \bigl(1+\ln \mathrm{tf}_{t,d}\bigr)\cdot
                \left(\ln\frac{1+N}{1+\mathrm{df}_t}+1\right),
  \label{eq:tfidf}
\end{equation}
after which every document vector is scaled to unit length, and the fifteen
engineered statistics are appended to the vector. Each base
learner is trained with imbalance handling, and its 5-fold cross-validated F1
on the training split determines its voice in the ensemble:
\begin{equation}
  p(\mathrm{spam}\mid x) \;=\; \sum_{m=1}^{3} w_m\, p_m(\mathrm{spam}\mid x),
  \qquad
  w_m \;=\; \frac{F1^{\mathrm{CV}}_m}{\sum_{k=1}^{3} F1^{\mathrm{CV}}_k},
  \label{eq:ensemble}
\end{equation}
so a stronger learner gets a proportionally louder vote. The held-out test
split is touched exactly once, for the final evaluation of Chapter~4.

\begin{figure}[H]
  \centering
  \begin{tikzpicture}[node distance=0.4cm and 1.0cm]
    \node[dat, minimum width=5.4cm] (a) {Download corpora (with mirror fallbacks)};
    \node[blk, below=of a, minimum width=5.4cm] (b) {Clean, de-duplicate, label};
    \node[blk, below=of b, minimum width=5.4cm] (c) {Stratified 80/20 train--test split};
    \node[blk, below=of c, minimum width=5.4cm] (d) {Normalise text $\to$ placeholder tokens};
    \node[blk, right=of a, minimum width=5.9cm] (e) {Featurise: TF--IDF + 15 stylometric};
    \node[blk, below=of e, minimum width=5.9cm] (f) {5-fold CV per base learner $\Rightarrow$ weights $w_m$};
    \node[blk, below=of f, minimum width=5.9cm] (g) {Fit LR, RF, XGBoost; build weighted ensemble};
    \node[blk, below=of g, minimum width=5.9cm] (h) {Evaluate once on held-out test set};
    \node[xai, below=of h, minimum width=5.9cm] (i) {Generate SHAP/LIME artefacts, figures, metrics};
    \draw[arr] (a) -- (b);
    \draw[arr] (b) -- (c);
    \draw[arr] (c) -- (d);
    \draw[arr] (d.east) .. controls +(0.9,0) and +(-0.9,0) .. (e.west);
    \draw[arr] (e) -- (f);
    \draw[arr] (f) -- (g);
    \draw[arr] (g) -- (h);
    \draw[arr] (h) -- (i);
  \end{tikzpicture}
  \caption{Implementation flowchart of the training and evaluation pipeline}
  \label{fig:flow}
\end{figure}

\subsection{Hyper-parameter configuration}

All hyper-parameters are kept in one configuration module
(\texttt{src/config.py}), so every experiment can be reproduced exactly.
Table~\ref{tab:hyper} lists the values used for both specialists. The learners
use standard, moderately regularised settings rather than aggressively tuned
ones, because the aim of the project is a transparent and reproducible system;
the three learners nevertheless reach near-identical cross-validated F1 scores
(Appendix-3).

\begin{table}[!htbp]
  \centering
  \caption{Hyper-parameters and experimental settings}
  \label{tab:hyper}
  \begin{tabular}{|L{3.4cm}|L{10.4cm}|}
    \hline
    \textbf{Component} & \textbf{Settings} \\
    \hline
    TF--IDF & SMS: 3,000 terms, minimum document frequency 2; e-mail: 5,000 terms, minimum document frequency 3; uni- and bi-grams; sublinear term frequency; accents stripped \\
    \hline
    Logistic Regression & \texttt{liblinear} solver; $C=1.0$; up to 2,000 iterations; balanced class weights \\
    \hline
    Random Forest & 250 trees; minimum of 2 samples per leaf; balanced class weights; all CPU cores \\
    \hline
    XGBoost & 300 trees; maximum depth 6; learning rate 0.1; row subsampling 0.9; column subsampling 0.8; histogram tree method; log-loss; $\mathrm{scale\_pos\_weight}=N_{\mathrm{ham}}/N_{\mathrm{spam}}$ on the training split (7.03 for SMS, 1.07 for e-mail) \\
    \hline
    Ensemble & Weights proportional to mean 5-fold cross-validated F1; decision threshold 0.5 \\
    \hline
    Protocol & Stratified 80/20 split; 5 cross-validation folds; random seed 42; e-mail sample of 12,000 messages \\
    \hline
    Explanations & LIME: 1,000 perturbations, 10 words per explanation; global SHAP: 800 held-out messages per modality \\
    \hline
  \end{tabular}
\end{table}

\subsection{Training and inference procedures}

The complete training procedure of one specialist is given as Algorithm~1.

\begin{center}
\fbox{\parbox{0.94\textwidth}{\small
\textbf{Algorithm 1: Training one channel specialist}\\[3pt]
\textbf{Input:} cleaned, de-duplicated corpus $D=\{(x_i,y_i)\}$ of one channel\\
\textbf{Output:} featuriser, three base learners, ensemble weights, evaluation report
\begin{enumerate}[leftmargin=0.9cm]
  \item Split $D$ into 80\% training and 20\% test data, stratified by label (seed 42).
  \item Fit the featuriser on the training texts only: learn the TF--IDF vocabulary of the normalised text and the maximum of each stylometric feature.
  \item Compute $\mathrm{scale\_pos\_weight}=N_{\mathrm{ham}}/N_{\mathrm{spam}}$ on the training split.
  \item For each learner $m\in\{\mathrm{LR},\mathrm{RF},\mathrm{XGB}\}$, estimate $\mathrm{F1}^{\mathrm{CV}}_m$ by stratified 5-fold cross-validation on the training split.
  \item Fit each learner on the full training split.
  \item Set $w_m=\mathrm{F1}^{\mathrm{CV}}_m/\sum_k \mathrm{F1}^{\mathrm{CV}}_k$ and form the weighted soft-voting ensemble.
  \item Evaluate the three learners and the ensemble once on the test split; save the metrics, the figures and the detector bundle.
\end{enumerate}}}
\end{center}

\subsection{Modality router}

The router decides which specialist scores a message when the user has not
chosen one. It is deliberately simple and fully transparent: a message is
treated as e-mail if it is longer than 280 characters, contains a line break,
or contains one of the words \emph{subject}, \emph{dear}, \emph{regards},
\emph{sincerely}, \emph{unsubscribe} or \emph{attached}; every other message
is treated as SMS. The 280-character limit lies well above the 160 characters
of a single SMS segment, and the cue words are typical of e-mail salutations,
sign-offs and mailing lists. Because a rule of this kind can be wrong, the
application lets the user override it, and its accuracy is measured in
Section~\ref{sec:routing}.

%------------------------------------------------------------------------------
\chapter{RESULTS ANALYSIS AND VALIDATION}

\section{Implementation of solution}

Modern engineering tools were used throughout:

\begin{itemize}
  \item \textbf{Analysis:} Python (pandas, NumPy) and a Jupyter research
        notebook for exploratory data analysis and sampling studies.
  \item \textbf{Design and implementation:} scikit-learn (Logistic
        Regression, Random Forest, TF--IDF), XGBoost, and the SHAP and LIME
        libraries, organised as a modular \texttt{src/} package.
  \item \textbf{Report preparation:} \LaTeX{} for this document; matplotlib
        with a shared publication style for all generated figures.
  \item \textbf{Project management and communication:} Git for version
        control and a shared Kanban board for task tracking.
  \item \textbf{Testing and validation:} a 13-check automated end-to-end
        smoke-test suite (including a simulated application session) and a
        strict held-out evaluation protocol.
\end{itemize}

Figure~\ref{fig:dash} shows the Detector tab of the Streamlit application
after a message has been analysed: the routing decision, the verdict with the
ensemble spam probability, and the vote of each base learner with its weight.

\begin{figure}[H]
  \centering
  \figimg[0.85\textwidth]{figures/app_detector_tab.png}{2.6in}
  \caption{Detector tab of the Streamlit application after analysing a message}
  \label{fig:dash}
\end{figure}

\subsection{Experimental setup}

Two public research corpora were used; their statistics after cleaning and
de-duplication are given in Table~\ref{tab:datasets} and visualised in
Figure~\ref{fig:classbalance}. Each corpus received an 80/20 stratified
train--test split; base-learner weights were computed by 5-fold
cross-validation on the training split only, and the test split was used
exactly once. Performance is reported with accuracy, precision
$P=\mathrm{TP}/(\mathrm{TP}+\mathrm{FP})$, recall
$R=\mathrm{TP}/(\mathrm{TP}+\mathrm{FN})$, their harmonic mean F1, and
ROC-AUC.

\begin{table}[H]
  \centering
  \caption{Corpora used for training and evaluation (after cleaning)}
  \label{tab:datasets}
  \begin{tabular}{|L{1.7cm}|L{5.6cm}|c|c|c|}
    \hline
    \textbf{Modality} & \textbf{Corpus} & \textbf{Messages} & \textbf{Spam share} & \textbf{Median words} \\
    \hline
    SMS & UCI SMS Spam Collection (Almeida \& G\'{o}mez Hidalgo, 2011) & 5,158 & 12.4\% & 12 \\
    \hline
    E-mail & Enron-Spam, stratified subsample (Metsis et al., 2006) & 12,000 & 48.3\% & 155 \\
    \hline
  \end{tabular}
\end{table}

\begin{figure}[H]
  \centering
  \figimg[0.9\textwidth]{figures/dataset_overview.png}{2.4in}
  \caption{Class balance per modality and message-length distributions}
  \label{fig:classbalance}
\end{figure}

The two modalities differ radically in both imbalance (12\% versus 48\%
spam) and length (12 versus 155 median words), which is precisely
why each is served by its own feature space and its own trained ensemble.

\textbf{Exploratory analysis of the vocabulary.} Before any model was trained,
the corpora were inspected for the words that best separate the two classes.
For every token present in at least eight messages after normalisation, the
log-ratio $\ln\bigl(P(\text{token}\mid\text{spam})/P(\text{token}\mid\text{ham})\bigr)$
of its presence probabilities was computed (positive values indicate spam). On SMS the strongest spam indicators are prize and
premium-rate vocabulary (\emph{claim}, \emph{prize}, \emph{guaranteed},
\emph{awarded}, \emph{ppm}, \emph{pobox}, \emph{ringtone}) together with the
placeholder \texttt{phonenumbr}, while ham is marked by informal chat
(\emph{lor}, \emph{lol}, \emph{cos}, \emph{later}) and by \emph{lt} and
\emph{gt}, the remains of the HTML-escaped placeholder \texttt{\&lt;\#\&gt;}
that replaces masked numbers in the corpus and appears in 197 legitimate
messages but in no spam. On e-mail the spam side is dominated by
pharmaceutical and software-piracy vocabulary (\emph{viagra}, \emph{cialis},
\emph{meds}, \emph{photoshop}, \emph{macromedia}), but the ham side is
dominated by Enron-specific tokens: the company name \emph{enron} itself
(log-ratio $-8.39$), the surname \emph{kaminski}, the abbreviations
\emph{hpl} and \emph{ena}, and energy-trading jargon (\emph{hourahead},
\emph{mmbtu}, \emph{westdesk}). This was the first sign of the dataset bias
that the explanation layer later confirmed (Section~\ref{sec:shortcut}).

\textbf{Statistical protocol.} Because all systems score above 97\%,
differences between them are small and must be checked for significance.
Uncertainty is quantified with 95\% percentile bootstrap confidence intervals
over 2,000 resamples of each test set (Efron and Tibshirani, 1993). The
ensemble is compared with each base learner on the same test messages with
the exact McNemar test on discordant predictions (McNemar, 1947) and with
DeLong's test for correlated ROC curves (DeLong et al., 1988); $p$-values are
corrected for the three comparisons per modality with Holm's method (Holm,
1979). The Matthews correlation coefficient (MCC), the area under the
precision--recall curve (PR-AUC), the false-positive rate (FPR) and the Brier
score of the predicted probabilities complement the headline metrics.

\subsection{Classification results}

Tables~\ref{tab:sms} and \ref{tab:email} report the complete held-out
results for all four systems on each modality; Figure~\ref{fig:metrics}
compares them graphically and Figure~\ref{fig:roc} shows the ROC curves.

\begin{table}[H]
  \centering
  \caption{SMS --- held-out test set (1,032 messages)}
  \label{tab:sms}
  \setlength{\tabcolsep}{4pt}
  \begin{tabular}{|L{3.3cm}|c|c|c|c|c|c|}
    \hline
    \textbf{Model} & \textbf{Accuracy} & \textbf{Precision} & \textbf{Recall} & \textbf{F1} & \textbf{ROC-AUC} & \textbf{CV F1} \\
    \hline
    Logistic Regression & 0.9855 & 0.9520 & 0.9297 & 0.9407 & 0.9941 & 0.9554 $\pm$ 0.020 \\
    \hline
    Random Forest & 0.9845 & 0.9516 & 0.9219 & 0.9365 & 0.9929 & 0.9558 $\pm$ 0.017 \\
    \hline
    XGBoost & 0.9845 & 0.9667 & 0.9062 & 0.9355 & 0.9941 & 0.9586 $\pm$ 0.015 \\
    \hline
    \textbf{Hybrid Ensemble} & \textbf{0.9855} & \textbf{0.9669} & \textbf{0.9141} & \textbf{0.9398} & \textbf{0.9941} & --- \\
    \hline
  \end{tabular}
\end{table}

\begin{table}[H]
  \centering
  \caption{E-mail --- held-out test set (2,400 messages)}
  \label{tab:email}
  \setlength{\tabcolsep}{4pt}
  \begin{tabular}{|L{3.3cm}|c|c|c|c|c|c|}
    \hline
    \textbf{Model} & \textbf{Accuracy} & \textbf{Precision} & \textbf{Recall} & \textbf{F1} & \textbf{ROC-AUC} & \textbf{CV F1} \\
    \hline
    Logistic Regression & 0.9833 & 0.9753 & 0.9905 & 0.9829 & 0.9979 & 0.9797 $\pm$ 0.004 \\
    \hline
    Random Forest & 0.9746 & 0.9629 & 0.9853 & 0.9740 & 0.9973 & 0.9669 $\pm$ 0.004 \\
    \hline
    XGBoost & 0.9750 & 0.9645 & 0.9845 & 0.9744 & 0.9969 & 0.9732 $\pm$ 0.003 \\
    \hline
    \textbf{Hybrid Ensemble} & \textbf{0.9825} & \textbf{0.9705} & \textbf{0.9940} & \textbf{0.9821} & \textbf{0.9983} & --- \\
    \hline
  \end{tabular}
\end{table}

\begin{figure}[H]
  \centering
  \begin{tikzpicture}
    \begin{axis}[
      ybar, bar width=4.2pt, width=0.44\textwidth, height=5.6cm,
      ymin=0.90, ymax=1.004,
      symbolic x coords={LR,RF,XGB,Ensemble}, xtick=data,
      enlarge x limits=0.16, ymajorgrids,
      title={SMS (test, $n=1032$)},
      tick label style={font=\scriptsize},
      label style={font=\scriptsize},
      title style={font=\footnotesize\bfseries},
      ylabel={Score},
      legend style={font=\scriptsize, at={(1.09,-0.24)}, anchor=north, legend columns=3},
    ]
      \addplot coordinates {(LR,0.9855)(RF,0.9845)(XGB,0.9845)(Ensemble,0.9855)};
      \addplot coordinates {(LR,0.9407)(RF,0.9365)(XGB,0.9355)(Ensemble,0.9398)};
      \addplot coordinates {(LR,0.9941)(RF,0.9929)(XGB,0.9941)(Ensemble,0.9941)};
      \legend{Accuracy,F1,ROC-AUC}
    \end{axis}
  \end{tikzpicture}\hfill
  \begin{tikzpicture}
    \begin{axis}[
      ybar, bar width=4.2pt, width=0.44\textwidth, height=5.6cm,
      ymin=0.90, ymax=1.004,
      symbolic x coords={LR,RF,XGB,Ensemble}, xtick=data,
      enlarge x limits=0.16, ymajorgrids,
      title={E-mail (test, $n=2400$)},
      tick label style={font=\scriptsize},
      label style={font=\scriptsize},
      title style={font=\footnotesize\bfseries},
    ]
      \addplot coordinates {(LR,0.9833)(RF,0.9746)(XGB,0.9750)(Ensemble,0.9825)};
      \addplot coordinates {(LR,0.9829)(RF,0.9740)(XGB,0.9744)(Ensemble,0.9821)};
      \addplot coordinates {(LR,0.9979)(RF,0.9973)(XGB,0.9969)(Ensemble,0.9983)};
    \end{axis}
  \end{tikzpicture}
  \caption{Test-set metrics per model and modality (vertical axis truncated at 0.90 to resolve the differences)}
  \label{fig:metrics}
\end{figure}

\begin{figure}[H]
  \centering
  \figimg[0.48\textwidth]{figures/sms_roc_curves.png}{2.3in}\hfill
  \figimg[0.48\textwidth]{figures/email_roc_curves.png}{2.3in}
  \caption{ROC curves for all four systems: SMS (left) and e-mail (right)}
  \label{fig:roc}
\end{figure}

On ROC-AUC the hybrid ensemble ties the best base learners on SMS (0.9941,
identical to four decimal places) and wins outright on e-mail (0.9983), where
it also posts near-perfect recall (0.9940) at 0.9705 precision. On SMS it
pairs the highest precision of any system (0.9669) with the joint-highest
accuracy (0.9855, tied with Logistic Regression) --- exactly the profile the
user-safety constraint asked for, since on that channel a false alarm means a
suppressed legitimate message.

Figure~\ref{fig:dotplots} shows the four threshold metrics side by side for
each modality. The plots make the scale of the differences visible: on SMS all
four systems lie within 0.001 of each other in accuracy, and the largest
spread, in recall (0.906--0.930), corresponds to only three of the 128 spam
messages in the test set.

\begin{figure}[!htbp]
  \centering
  \figimg[0.82\textwidth]{figures/sms_model_comparison.png}{2.4in}\\[6pt]
  \figimg[0.82\textwidth]{figures/email_model_comparison.png}{2.4in}
  \caption{Test-set accuracy, precision, recall and F1 of the four systems: SMS (top) and e-mail (bottom)}
  \label{fig:dotplots}
\end{figure}

\textbf{Additional metrics.} Table~\ref{tab:extrametrics} reports metrics that
are more informative than accuracy under class imbalance. The ensemble has the
best PR-AUC on both modalities (0.9790 on SMS, 0.9980 on e-mail) and the best
Brier score on SMS (0.0130), so its probabilities are both well ranked and
accurate there; on e-mail, XGBoost and Logistic Regression have slightly lower
Brier scores. Its MCC (0.9320 on SMS, 0.9652 on e-mail) is within 0.002 of the
best learner on both modalities, and its SMS false-positive rate of 0.44\%
equals the best, that of XGBoost.

\begin{table}[!htbp]
  \centering
  \caption{Additional metrics on the held-out test sets (lower Brier score and FPR are better)}
  \label{tab:extrametrics}
  \begin{tabular}{|L{1.8cm}|L{3.8cm}|c|c|c|c|}
    \hline
    \textbf{Modality} & \textbf{Model} & \textbf{MCC} & \textbf{PR-AUC} & \textbf{Brier} & \textbf{FPR} \\
    \hline
    \multirow{4}{*}{SMS} & Logistic Regression & 0.9325 & 0.9768 & 0.0148 & 0.0066 \\
    \cline{2-6}
     & Random Forest & 0.9278 & 0.9770 & 0.0184 & 0.0066 \\
    \cline{2-6}
     & XGBoost & 0.9273 & 0.9739 & 0.0135 & 0.0044 \\
    \cline{2-6}
     & \textbf{Hybrid Ensemble} & 0.9320 & 0.9790 & 0.0130 & 0.0044 \\
    \hline
    \multirow{4}{*}{E-mail} & Logistic Regression & 0.9668 & 0.9976 & 0.0210 & 0.0233 \\
    \cline{2-6}
     & Random Forest & 0.9494 & 0.9970 & 0.0380 & 0.0354 \\
    \cline{2-6}
     & XGBoost & 0.9502 & 0.9965 & 0.0189 & 0.0338 \\
    \cline{2-6}
     & \textbf{Hybrid Ensemble} & 0.9652 & 0.9980 & 0.0213 & 0.0282 \\
    \hline
  \end{tabular}
\end{table}

\textbf{Confidence intervals.} Table~\ref{tab:ci} gives the 95\% bootstrap
confidence intervals of the ensemble. They are narrow on e-mail (about
$\pm$0.5 percentage points of accuracy) but wide on SMS for the spam-class
metrics: with only 128 spam messages in the SMS test set, the recall interval
spans 0.861--0.960. Differences of a few thousandths between systems on SMS
are therefore well inside the sampling uncertainty.

\begin{table}[!htbp]
  \centering
  \caption{Hybrid ensemble: point estimates and 95\% bootstrap confidence intervals (2,000 resamples of the test set)}
  \label{tab:ci}
  \begin{tabular}{|L{2.6cm}|c|c|c|c|}
    \hline
    \textbf{Metric} & \textbf{SMS} & \textbf{95\% CI} & \textbf{E-mail} & \textbf{95\% CI} \\
    \hline
    Accuracy & 0.9855 & [0.9777, 0.9922] & 0.9825 & [0.9771, 0.9875] \\
    \hline
    Precision & 0.9669 & [0.9286, 0.9925] & 0.9705 & [0.9607, 0.9799] \\
    \hline
    Recall & 0.9141 & [0.8605, 0.9600] & 0.9940 & [0.9894, 0.9982] \\
    \hline
    F1 & 0.9398 & [0.9068, 0.9669] & 0.9821 & [0.9766, 0.9873] \\
    \hline
    MCC & 0.9320 & [0.8954, 0.9625] & 0.9652 & [0.9547, 0.9751] \\
    \hline
    ROC-AUC & 0.9941 & [0.9882, 0.9987] & 0.9983 & [0.9974, 0.9990] \\
    \hline
    PR-AUC & 0.9790 & [0.9610, 0.9934] & 0.9980 & [0.9971, 0.9989] \\
    \hline
  \end{tabular}
\end{table}

\subsection{Statistical significance of the differences}

Table~\ref{tab:signif} lists the paired tests. On SMS none of the differences
is significant: the ensemble and each base learner disagree on at most eight
of the 1,032 test messages, and the ROC-AUC differences are below 0.0013. On
e-mail the picture is clearer. The ensemble is significantly better than both
tree models --- it classifies 24 messages correctly that Random Forest gets
wrong, against 5 in the other direction, and 22 against 4 for XGBoost
(Holm-corrected McNemar $p=0.002$) --- and its ROC-AUC is significantly higher
(DeLong $p\le0.003$). Against Logistic Regression, however, the counts are
balanced (11 against 13) and neither test is significant. Combining a linear
model with two tree ensembles therefore lifts the tree models to the level of
Logistic Regression and improves the ranking metrics, but on these corpora it
does not surpass a well-regularised linear model on TF--IDF features.

\begin{table}[!htbp]
  \centering
  \caption{Paired significance tests of the hybrid ensemble against each base learner ($b$ / $c$: messages that only the ensemble / only the base learner classifies correctly; $p$-values Holm-corrected)}
  \label{tab:signif}
  \setlength{\tabcolsep}{4pt}
  \begin{tabular}{|L{1.8cm}|L{3.6cm}|c|c|c|c|}
    \hline
    \textbf{Modality} & \textbf{Comparison} & $b$\,/\,$c$ & \textbf{McNemar} $p$ & $\Delta$\textbf{AUC} & \textbf{DeLong} $p$ \\
    \hline
    SMS & Ensemble vs.\ Logistic Regression & 4 / 4 & 1.000 & +0.0000 & 1.000 \\
    \hline
    SMS & Ensemble vs.\ Random Forest & 2 / 1 & 1.000 & +0.0012 & 0.754 \\
    \hline
    SMS & Ensemble vs.\ XGBoost & 1 / 0 & 1.000 & -0.0000 & 1.000 \\
    \hline
    E-mail & Ensemble vs.\ Logistic Regression & 11 / 13 & 0.839 & +0.0004 & 0.208 \\
    \hline
    E-mail & Ensemble vs.\ Random Forest & 24 / 5 & 0.002 & +0.0009 & 0.003 \\
    \hline
    E-mail & Ensemble vs.\ XGBoost & 22 / 4 & 0.002 & +0.0014 & $<$0.001 \\
    \hline
  \end{tabular}
\end{table}

\subsection{Ensemble behaviour and error analysis}

Figure~\ref{fig:confusion} shows the ensemble confusion matrices derived from
the held-out evaluations. The SMS ensemble mislabels only 4 of 904 legitimate
messages; the e-mail ensemble misses just 7 of 1,158 spam messages. The two
matrices are deliberately asymmetric: the SMS specialist errs on the side of
precision, the e-mail specialist on the side of recall, matching the relative
costs of a wrong decision on each channel.

\begin{figure}[H]
  \centering
  \begin{minipage}{0.48\textwidth}\centering
  \begin{tikzpicture}[font=\footnotesize]
    \def\cw{2.15}\def\ch{1.2}
    \fill[green!14] (0,\ch) rectangle (\cw,2*\ch);
    \fill[red!12]   (\cw,\ch) rectangle (2*\cw,2*\ch);
    \fill[red!12]   (0,0) rectangle (\cw,\ch);
    \fill[green!14] (\cw,0) rectangle (2*\cw,\ch);
    \draw (0,0) rectangle (2*\cw,2*\ch);
    \draw (\cw,0) -- (\cw,2*\ch);
    \draw (0,\ch) -- (2*\cw,\ch);
    \node[align=center] at (0.5*\cw,1.5*\ch) {TN\\\textbf{900}};
    \node[align=center] at (1.5*\cw,1.5*\ch) {FP\\\textbf{4}};
    \node[align=center] at (0.5*\cw,0.5*\ch) {FN\\\textbf{11}};
    \node[align=center] at (1.5*\cw,0.5*\ch) {TP\\\textbf{117}};
    \node[anchor=south] at (0.5*\cw,2*\ch+0.05) {Pred.\ ham};
    \node[anchor=south] at (1.5*\cw,2*\ch+0.05) {Pred.\ spam};
    \node[anchor=east]  at (-0.06,1.5*\ch) {Actual ham};
    \node[anchor=east]  at (-0.06,0.5*\ch) {Actual spam};
    \node[anchor=south, font=\footnotesize\bfseries] at (\cw,2*\ch+0.5) {SMS ensemble ($n=1032$)};
  \end{tikzpicture}
  \end{minipage}\hfill
  \begin{minipage}{0.48\textwidth}\centering
  \begin{tikzpicture}[font=\footnotesize]
    \def\cw{2.15}\def\ch{1.2}
    \fill[green!14] (0,\ch) rectangle (\cw,2*\ch);
    \fill[red!12]   (\cw,\ch) rectangle (2*\cw,2*\ch);
    \fill[red!12]   (0,0) rectangle (\cw,\ch);
    \fill[green!14] (\cw,0) rectangle (2*\cw,\ch);
    \draw (0,0) rectangle (2*\cw,2*\ch);
    \draw (\cw,0) -- (\cw,2*\ch);
    \draw (0,\ch) -- (2*\cw,\ch);
    \node[align=center] at (0.5*\cw,1.5*\ch) {TN\\\textbf{1207}};
    \node[align=center] at (1.5*\cw,1.5*\ch) {FP\\\textbf{35}};
    \node[align=center] at (0.5*\cw,0.5*\ch) {FN\\\textbf{7}};
    \node[align=center] at (1.5*\cw,0.5*\ch) {TP\\\textbf{1151}};
    \node[anchor=south] at (0.5*\cw,2*\ch+0.05) {Pred.\ ham};
    \node[anchor=south] at (1.5*\cw,2*\ch+0.05) {Pred.\ spam};
    \node[anchor=east]  at (-0.06,1.5*\ch) {Actual ham};
    \node[anchor=east]  at (-0.06,0.5*\ch) {Actual spam};
    \node[anchor=south, font=\footnotesize\bfseries] at (\cw,2*\ch+0.5) {E-mail ensemble ($n=2400$)};
  \end{tikzpicture}
  \end{minipage}
  \caption{Confusion matrices of the hybrid ensembles on the held-out test sets}
  \label{fig:confusion}
\end{figure}

The confusion matrices of all four systems (Figure~\ref{fig:confall}) show
where the ensemble's advantage comes from. On e-mail it has the lowest
false-negative count of all systems (7, against 11 for Logistic Regression, 17
for Random Forest and 18 for XGBoost), with a false-positive count (35)
between that of Logistic Regression (29) and the tree models (42 and 44). On
SMS it matches the four false positives of XGBoost while missing one spam
message fewer (11 against 12).

\begin{figure}[!htbp]
  \centering
  \figimg[0.49\textwidth]{figures/sms_confusion_matrices.png}{2.4in}\hfill
  \figimg[0.49\textwidth]{figures/email_confusion_matrices.png}{2.4in}
  \caption{Confusion matrices of all four systems on the held-out test sets: SMS (left) and e-mail (right)}
  \label{fig:confall}
\end{figure}

\textbf{Error analysis.} Table~\ref{tab:errors} lists representative
misclassified test messages with the ensemble's spam probability. The errors
are instructive. The SMS false positives share surface traits of spam: prices
and a deposit, informal talk about money, and a
genuine bank security notice whose vocabulary (\emph{sensitive information},
\emph{passwords}, \emph{PIN}) is exactly that of phishing. The SMS false
negatives are mostly joke or quiz messages written in a
conversational style, with few of the digits, prices and links on which the
stylometric features rely. On e-mail, the false positives include short
forwarded notes and external newsletters, and none of the 35 contains the
token \emph{enron}; the missed spam includes an advance-fee fraud written in
plain business language and a one-line lure to open an attachment, both of
which lack the vocabulary that characterises most of the spam class.

\begin{table}[!htbp]
  \centering
  \caption{Representative misclassified test messages (ensemble spam probability; threshold 0.5)}
  \label{tab:errors}
  \small
  \begin{tabular}{|L{1.7cm}|L{10.3cm}|c|}
    \hline
    \textbf{Error} & \textbf{Message (truncated)} & \textbf{$p$(spam)} \\
    \hline
    SMS FP & Please protect yourself from e-threats. SIB never asks for sensitive information like Passwords,ATM/SMS PIN \ldots & 0.816 \\
    \hline
    SMS FP & Hey\ldots Great deal\ldots Farm tour 9am to 5pm \$95/pax, \$50 deposit by 16 May & 0.778 \\
    \hline
    SMS FN & Latest News! Police station toilet stolen, cops have nothing to go on! & 0.065 \\
    \hline
    SMS FN & In The Simpsons Movie released in July 2007 name the band that died at the start of the film? A-Green Day, \ldots & 0.202 \\
    \hline
    E-mail FP & the negotiation skills company , inc . newsletter aug 2001 fighting fires without burning bridges \ldots & 0.824 \\
    \hline
    E-mail FP & fw : very symbolic of our current struggle & 0.675 \\
    \hline
    E-mail FN & please read urgently dear sir / madam , i am contacting you from asia to discuss a business deal of us \$ 18 \ldots & 0.308 \\
    \hline
    E-mail FN & re : your details please read the attached file . & 0.323 \\
    \hline
  \end{tabular}
\end{table}

\subsection{Routing and cross-channel transfer}\label{sec:routing}

The per-modality results assume that every message reaches the correct
specialist. Table~\ref{tab:routing} tests that assumption on the pooled test
sets (3,432 messages). The specialists do not transfer: the SMS model is no
better than chance on e-mail (ROC-AUC 0.53, MCC 0.04), and the e-mail model
labels 88\% of legitimate SMS as spam, so channel-specific training is
essential. The heuristic router sends 96.4\% of SMS and 85.1\% of e-mails to
the right specialist; the misrouted e-mails are short and contain none of the
cue words, and the SMS specialist classifies only 42.3\% of them correctly. As
a result, end-to-end accuracy on the pooled test sets falls from 98.34\% with
oracle routing to 92.16\% with automatic routing. The classifiers are
accurate; the router is the weakest component of the deployed system, which
is why the application lets the user override it. Figure~\ref{fig:dash} shows such a case: an SMS containing the word \emph{unsubscribe} is
routed to the e-mail ensemble.

\begin{table}[!htbp]
  \centering
  \caption{Routing and cross-channel transfer of the ensembles (pooled test set of 3,432 messages)}
  \label{tab:routing}
  \setlength{\tabcolsep}{4pt}
  \begin{tabular}{|L{5.6cm}|c|c|c|c|c|}
    \hline
    \textbf{Setting} & \textbf{Accuracy} & \textbf{F1} & \textbf{MCC} & \textbf{AUC} & \textbf{FPR} \\
    \hline
    Oracle routing (true channel), pooled & 0.9834 & 0.9780 & 0.9648 & 0.9978 & 0.0182 \\
    \hline
    Automatic routing, pooled & 0.9216 & 0.8921 & 0.8317 & 0.9687 & 0.0443 \\
    \hline
    SMS model applied to e-mail test set & 0.5208 & 0.5287 & 0.0442 & 0.5290 & 0.5129 \\
    \hline
    E-mail model applied to SMS test set & 0.2209 & 0.2299 & 0.0594 & 0.6707 & 0.8805 \\
    \hline
  \end{tabular}
\end{table}

\subsection{Explainability results}

\subsubsection{A verdict, explained (local scope)}

The following scam SMS, never seen in training, is used as a live exhibit:

\begin{center}
\fbox{\parbox{0.88\textwidth}{\small
``URGENT! You have WON a \pounds 2,000 cash prize GUARANTEED. Call
09061701461 now or visit www.claim-your-prize.com to claim!''}}
\end{center}

All three base learners vote spam --- Logistic Regression 1.000, Random
Forest 0.987, XGBoost 1.000 --- giving a weighted ensemble probability of
\textbf{0.996 spam}. LIME re-scores approximately 1,000 word-dropout
perturbations of the message through the entire pipeline and fits a local
linear surrogate; the resulting per-word contributions are plotted in
Figure~\ref{fig:lime} alongside the SHAP waterfall of the same prediction.
The premium-rate phone number is the single strongest signal, followed by the
money amount and the classic scam vocabulary --- a justification a non-expert
can check at a glance.

\begin{figure}[H]
  \centering
  \begin{minipage}[c]{0.5\textwidth}\centering
  \begin{tikzpicture}
    \begin{axis}[
      xbar, width=0.76\textwidth, height=6.4cm,
      xmin=0, xmax=0.15,
      symbolic y coords={your,www,WON,GUARANTEED,2,prize,URGENT,claim,000,09061701461},
      ytick=data,
      yticklabel style={font=\scriptsize\ttfamily},
      tick label style={font=\scriptsize},
      xlabel={LIME weight toward ``spam''},
      label style={font=\scriptsize},
      nodes near coords, nodes near coords style={font=\tiny},
      /pgf/number format/.cd, fixed, precision=3,
      bar width=5.5pt,
    ]
      \addplot coordinates {
        (0.025,your) (0.031,www) (0.036,WON) (0.037,GUARANTEED) (0.038,2)
        (0.045,prize) (0.046,URGENT) (0.053,claim) (0.061,000) (0.123,09061701461)};
    \end{axis}
  \end{tikzpicture}
  \end{minipage}\hfill
  \begin{minipage}[c]{0.46\textwidth}\centering
  \figimg[0.98\textwidth]{figures/sms_shap_waterfall.png}{2.4in}
  \end{minipage}
  \caption{Local explanations of the exhibited SMS: LIME word attributions over the full ensemble (left) and SHAP waterfall of the XGBoost component (right)}
  \label{fig:lime}
\end{figure}

Figure~\ref{fig:limecases} applies LIME to two real spam messages from the
test sets, each the highest-scoring among the first 200 spam messages of its
test set. For the SMS, the premium-rate number, the PO-box code and
\emph{URGENT} carry most of the weight; for the e-mail, the link and
promotional words (\emph{online}, \emph{save}, \emph{free}) push towards spam,
while \emph{thanks} --- typical of legitimate correspondence --- pushes the
other way, showing that LIME reports evidence in both directions.

\begin{figure}[!htbp]
  \centering
  \figimg[0.49\textwidth]{figures/sms_lime_case_study.png}{2.2in}\hfill
  \figimg[0.49\textwidth]{figures/email_lime_case_study.png}{2.2in}
  \caption{LIME explanations of two held-out spam messages over the full ensemble: SMS (left) and e-mail (right); red bars push towards spam, blue towards ham}
  \label{fig:limecases}
\end{figure}

\subsubsection{Are the explanations faithful?}

An explanation is useful only if it reflects what the model actually relies
on. This was tested with a deletion experiment: for 60 SMS and 25 e-mail spam
messages that the ensemble detected correctly, the top-$k$ words ranked by
LIME ($k=1,3,5$) were removed and the drop in the ensemble's spam probability
was recorded, and compared with removing $k$ randomly chosen words (averaged
over ten draws). Table~\ref{tab:limefaith} shows that LIME's words matter far
more than random ones. On SMS, removing the single top word lowers the spam
probability by 0.106 on average, 8.6 times more than a random word, and
removing the top five lowers it by 0.455 and flips 52\% of the verdicts to
ham. On e-mail the drops are smaller in absolute terms, because long messages
contain redundant evidence, but still about 20 times larger than for random
deletions.

\begin{table}[!htbp]
  \centering
  \caption{Deletion test of LIME faithfulness: mean drop in the ensemble spam probability after removing the top-$k$ LIME words / $k$ random words}
  \label{tab:limefaith}
  \begin{tabular}{|L{3.0cm}|c|c|c|c|}
    \hline
    \textbf{Modality} & $k=1$ & $k=3$ & $k=5$ & \textbf{Flipped} ($k=5$) \\
    \hline
    SMS ($n=60$) & 0.106 / 0.012 & 0.291 / 0.043 & 0.455 / 0.081 & 52\% \\
    \hline
    E-mail ($n=25$) & 0.043 / 0.002 & 0.133 / 0.005 & 0.198 / 0.010 & 24\% \\
    \hline
  \end{tabular}
\end{table}

\subsubsection{What the models learned (global scope)}

Global SHAP summaries were computed over 800 test messages per modality
(Figure~\ref{fig:beeswarm}). The key finding is that \emph{stylometry rivals
vocabulary}: on SMS the five top-ranked features are all engineered
stylometric statistics --- \texttt{DIGIT\_RATIO}, \texttt{AVG\_WORD\_LEN},
\texttt{SPAM\_KEYWORDS}, \texttt{UPPERCASE\_RATIO} and \texttt{CHAR\_COUNT}
--- and the fifteen engineered features together carry about three-quarters
(74.5\%) of the total attribution, although they are only 15 of the 3,015
features. A scam is betrayed by its \emph{shape} as much as by its words.

All SHAP values in this chapter are computed with XGBoost's native, exact
TreeSHAP on the same sparse feature matrix that the model scores. This
matters: XGBoost treats entries that are absent from a sparse matrix as
missing values, so explaining a densified copy of the input would describe a
different function --- in this pipeline, densification changes the XGBoost
verdict for 729 of the 1,032 SMS test messages.

\begin{figure}[H]
  \centering
  \figimg[0.48\textwidth]{figures/sms_shap_beeswarm.png}{2.6in}\hfill
  \figimg[0.48\textwidth]{figures/email_shap_beeswarm.png}{2.6in}
  \caption{Global SHAP beeswarm summaries: SMS (left) and e-mail (right); UPPER-CASE rows are engineered stylometric features, lower-case rows are words}
  \label{fig:beeswarm}
\end{figure}

Figure~\ref{fig:shapimp} condenses the beeswarms into the mean absolute
attribution of each feature, and Tables~\ref{tab:shapsms}
and~\ref{tab:shapemail} list the ten most important features with the
direction of their effect. On SMS, eight of the ten are stylometric: a high
share of digits, longer words, scam-lexicon hits, capital letters and a longer
message all push towards spam, while a high share of punctuation and special
characters, typical of informal chat, pushes towards ham. On e-mail the
ranking is led by a single word, \emph{enron}, whose mean attribution (1.327)
is three times that of the next feature; the first name \emph{vince} and
correspondence words such as \emph{thanks} and \emph{attached} also push
towards ham, whereas exclamation marks, links and the words \emph{your} and
\emph{our} push towards spam.

\begin{figure}[!htbp]
  \centering
  \figimg[0.49\textwidth]{figures/sms_shap_importance.png}{2.3in}\hfill
  \figimg[0.49\textwidth]{figures/email_shap_importance.png}{2.3in}
  \caption{Global feature importance (mean $|$SHAP$|$, exact TreeSHAP on the XGBoost component, 800 held-out messages): SMS (left) and e-mail (right)}
  \label{fig:shapimp}
\end{figure}

\begin{table}[!htbp]
  \centering
  \caption{SMS: the ten most important features by mean $|$SHAP$|$ (exact TreeSHAP, XGBoost component, 800 held-out messages); the last column shows where high values of the feature push the prediction}
  \label{tab:shapsms}
  \begin{tabular}{|c|L{4.2cm}|L{3.4cm}|c|L{2.4cm}|}
    \hline
    \textbf{Rank} & \textbf{Feature} & \textbf{Type} & \textbf{Mean $|\phi|$} & \textbf{Pushes towards} \\
    \hline
    1 & \texttt{DIGIT\_RATIO} & Stylometric & 3.539 & spam \\
    \hline
    2 & \texttt{AVG\_WORD\_LEN} & Stylometric & 1.414 & spam \\
    \hline
    3 & \texttt{SPAM\_KEYWORDS} & Stylometric & 1.285 & spam \\
    \hline
    4 & \texttt{UPPERCASE\_RATIO} & Stylometric & 0.867 & spam \\
    \hline
    5 & \texttt{CHAR\_COUNT} & Stylometric & 0.851 & spam \\
    \hline
    6 & \texttt{SPECIAL\_CHAR\_RATIO} & Stylometric & 0.532 & ham \\
    \hline
    7 & \texttt{you} & Word / placeholder & 0.523 & spam \\
    \hline
    8 & \texttt{PUNCT\_RATIO} & Stylometric & 0.457 & ham \\
    \hline
    9 & \texttt{to} & Word / placeholder & 0.379 & spam \\
    \hline
    10 & \texttt{NUM\_EXCLAMATIONS} & Stylometric & 0.368 & spam \\
    \hline
  \end{tabular}
\end{table}

\begin{table}[!htbp]
  \centering
  \caption{E-mail: the ten most important features by mean $|$SHAP$|$ (exact TreeSHAP, XGBoost component, 800 held-out messages); the last column shows where high values of the feature push the prediction}
  \label{tab:shapemail}
  \begin{tabular}{|c|L{4.2cm}|L{3.4cm}|c|L{2.4cm}|}
    \hline
    \textbf{Rank} & \textbf{Feature} & \textbf{Type} & \textbf{Mean $|\phi|$} & \textbf{Pushes towards} \\
    \hline
    1 & \texttt{enron} & Word / placeholder & 1.327 & ham \\
    \hline
    2 & \texttt{NUM\_EXCLAMATIONS} & Stylometric & 0.432 & spam \\
    \hline
    3 & \texttt{CHAR\_COUNT} & Stylometric & 0.318 & spam \\
    \hline
    4 & \texttt{your} & Word / placeholder & 0.281 & spam \\
    \hline
    5 & \texttt{thanks} & Word / placeholder & 0.260 & ham \\
    \hline
    6 & \texttt{vince} & Word / placeholder & 0.239 & ham \\
    \hline
    7 & \texttt{http} & Word / placeholder & 0.238 & spam \\
    \hline
    8 & \texttt{attached} & Word / placeholder & 0.236 & ham \\
    \hline
    9 & \texttt{DIGIT\_RATIO} & Stylometric & 0.229 & ham \\
    \hline
    10 & \texttt{our} & Word / placeholder & 0.213 & spam \\
    \hline
  \end{tabular}
\end{table}

\subsubsection{XAI as a debugging tool: a dataset bias, caught}\label{sec:shortcut}

On the e-mail model, the global SHAP ranking exposed the token
\texttt{enron} as a strong \emph{ham} indicator: the corpus's legitimate mail
comes entirely from the mailboxes of Enron employees, so the model partly
learned the dataset rather than the concept. This bias is invisible in every accuracy table in this chapter
and obvious in one beeswarm plot --- the clearest argument this project makes
for shipping explanations together with predictions.

\textbf{Quantifying the shortcut.} The token occurs in 44.4\% of the
legitimate training e-mails and 42.1\% of the legitimate test e-mails, but in
none of the spam messages in either split, so it is a perfect one-sided
signal for ham. The exact attributions confirm that the model relies on it:
its mean TreeSHAP contribution is $-3.95$ log-odds when the token is present,
against $+0.66$ when it is absent. Two probes measured the practical
consequences (Table~\ref{tab:shortcut}). Deleting the token from the 523
legitimate test e-mails that contain it flips only 11 verdicts, because
legitimate mail contains other ham evidence. Appending the single word
\emph{enron} to spam, however, flips 409 of the 1,151 spam messages the
ensemble had detected (35.5\%) to ham. The learners differ sharply: XGBoost is
evaded on 91.7\% of the messages, Random Forest on 25.3\% and Logistic
Regression on only 4.5\%. The heterogeneous ensemble is therefore far more
robust than its boosted member, though not immune --- a concrete argument for
combining learners with different inductive biases.

\begin{table}[!htbp]
  \centering
  \caption{Shortcut probe on the e-mail test set: appending the word ``enron'' to detected spam, and deleting it from legitimate mail ($\overline{\Delta p}$: mean change in spam probability)}
  \label{tab:shortcut}
  \setlength{\tabcolsep}{4pt}
  \begin{tabular}{|L{3.6cm}|c|c|c|c|}
    \hline
    \textbf{Model} & \textbf{Spam flipped to ham} & $\overline{\Delta p}$ & \textbf{Ham flipped to spam} & $\overline{\Delta p}$ \\
    \hline
    Logistic Regression & 52 / 1147 (4.5\%) & -0.066 & 2 / 523 & +0.032 \\
    \hline
    Random Forest & 289 / 1141 (25.3\%) & -0.301 & 10 / 523 & +0.053 \\
    \hline
    XGBoost & 1045 / 1140 (91.7\%) & -0.763 & 19 / 523 & +0.052 \\
    \hline
    \textbf{Hybrid Ensemble} & 409 / 1151 (35.5\%) & -0.376 & 11 / 523 & +0.046 \\
    \hline
  \end{tabular}
\end{table}

\subsection{Computational cost}\label{sec:cost}

Table~\ref{tab:cost} summarises the cost of training and of explanation.
Training a specialist, including 5-fold cross-validation of all three
learners, takes 14 seconds for SMS and 77 seconds for e-mail on a CPU. A
single verdict takes about 15\,ms, an exact TreeSHAP attribution under a
millisecond and a LIME explanation with 1,000 perturbations 0.08--0.21\,s, so
every prediction can be explained while the user waits. Batched scoring
processes about 15,000 SMS or 2,000 e-mails per second.

\begin{table}[!htbp]
  \centering
  \caption{Computational cost on a 10-core ARM64 CPU without a GPU (median times)}
  \label{tab:cost}
  \setlength{\tabcolsep}{4pt}
  \begin{tabular}{|L{1.8cm}|c|c|c|c|c|}
    \hline
    \textbf{Modality} & \textbf{Training (s)} & \textbf{Verdict (ms)} & \textbf{Batch (msg/s)} & \textbf{SHAP (ms)} & \textbf{LIME (ms)} \\
    \hline
    SMS & 14.3 & 14.7 & 14{,}897 & 0.86 & 76 \\
    \hline
    E-mail & 76.8 & 14.8 & 1{,}967 & 0.75 & 210 \\
    \hline
  \end{tabular}
\end{table}

\subsection{Validation on realistic unseen messages}

The public corpora are more than a decade old, so the system was also tested
on 28 realistic messages written for this project in the style of current
scams and everyday messages: lottery and cash-back offers, electricity
disconnection threats, parcel and KYC scams, part-time job offers, account
suspension and password-expiry phishing, alongside ordinary personal texts,
meeting notes and university e-mails. None of these messages appears in
either corpus. With automatic routing, all 28 reached the intended specialist
and 26 were classified correctly; Table~\ref{tab:demo} lists a representative
selection together with both errors. The first error is a genuine one-time
password that receives a spam probability of 52.2\%: OTP messages are full of
digits, exactly like scams, which is why the SMS specialist was tuned for
precision. The second is an invoice-fraud e-mail (38.9\%) that contains none
of the vocabulary typical of spam; Logistic Regression flags it (65\%) but the
tree models do not.

The shortcut of Section~\ref{sec:shortcut} can also be demonstrated live. A
password-expiry phishing e-mail receives a spam probability of 86.6\%; signing
the same message as ``IT Support, Enron North America'' lowers it to 42.3\%
and flips the verdict to ham, with XGBoost dropping from 86.8\% to 3.1\% while
Logistic Regression still reports 80.4\%.

{\small
\begin{longtable}{|L{1.4cm}|L{7.4cm}|L{1.1cm}|L{2.3cm}|L{1.5cm}|}
  \caption{Validation on realistic messages written for this project (not taken from either corpus); spam probability from the routed ensemble}
  \label{tab:demo}\\
  \hline
  \textbf{Channel} & \textbf{Message} & \textbf{True} & \textbf{Predicted} & \textbf{Correct} \\
  \hline
  \endfirsthead
  \hline
  \textbf{Channel} & \textbf{Message} & \textbf{True} & \textbf{Predicted} & \textbf{Correct} \\
  \hline
  \endhead
  SMS & Congratulations! Your mobile number has won Rs 25,00,000 in our annual lucky draw. To claim your prize call 9876543210 now and share your bank details. & spam & spam (99.0\%) & Yes \\
  \hline
  SMS & Your electricity connection will be disconnected tonight at 9:30 PM due to an unpaid bill. Call our officer on 9123456780 immediately to avoid disconnection. & spam & spam (91.6\%) & Yes \\
  \hline
  SMS & Hey, are we still meeting for lunch at 1pm? I might be 10 min late. & ham & ham (7.9\%) & Yes \\
  \hline
  SMS & Mom, I reached the hostel safely. Will call you after dinner. & ham & ham (3.4\%) & Yes \\
  \hline
  E-mail & Subject: Password expiry notice. Your mailbox password expires today. Click http://mail-verify-portal.com to keep your account active and avoid losing your messages. IT Support & spam & spam (86.6\%) & Yes \\
  \hline
  E-mail & Subject: Your account has been suspended. Dear customer, we detected unusual activity. Verify your account details at http://secure-account-check.com within 24 hours or your account will be closed. Customer Service & spam & spam (85.2\%) & Yes \\
  \hline
  E-mail & Subject: Q3 planning meeting notes. Hi team, attached are the notes from yesterday's planning session. Please review the action items assigned to you before Friday's follow-up. Thanks, Sarah & ham & ham (14.0\%) & Yes \\
  \hline
  SMS & Your OTP for login is 482913. It is valid for 10 minutes. Do not share it with anyone. & ham & spam (52.2\%) & \textbf{No} \\
  \hline
  E-mail & Subject: Invoice payment overdue. Please find the attached invoice and make the payment of 2,450 dollars today using the link http://pay-invoice-online.net to avoid late fees. Accounts Department & spam & ham (38.9\%) & \textbf{No} \\
  \hline
\end{longtable}
}

\subsection{Interpretation and data validation}

Cross-validation dispersion is small (F1 standard deviation between
$\pm$0.003 and $\pm$0.020 across all six base learners), indicating stable
models rather than lucky splits, and the derived ensemble weights are
near-uniform (Appendix-3), so the ensemble's gains come from the
\emph{diversity} of the three learners rather than from one dominant model.
The 13-check automated smoke-test suite exercises the entire system ---
data loading, preprocessing, both detectors, LIME, SHAP and a simulated
application session --- and passes end-to-end; the complete pipeline
reproduces from a clean checkout in about three minutes on a commodity CPU,
with training itself taking about two minutes. The dataset downloader
validates its sources and falls back to mirrors automatically, and the
modality router can be overridden by the user in the application, so a
misrouted message can always be re-scored by the other specialist.

Table~\ref{tab:smoke} lists the thirteen automated checks and the result of
the final run.

\begin{table}[!htbp]
  \centering
  \caption{Automated end-to-end smoke tests (final run: all checks passed)}
  \label{tab:smoke}
  \begin{tabular}{|c|L{8.6cm}|L{4.2cm}|}
    \hline
    \textbf{No.} & \textbf{Check} & \textbf{Result} \\
    \hline
    1 & SMS detector flags a scam SMS as spam & Pass ($p=0.996$) \\
    \hline
    2 & SMS detector passes an ordinary SMS as ham & Pass ($p=0.079$) \\
    \hline
    3 & All three SMS base learners return a vote & Pass \\
    \hline
    4 & LIME returns word weights for the scam SMS & Pass \\
    \hline
    5 & SHAP returns attributions for the scam SMS & Pass \\
    \hline
    6 & E-mail detector flags a phishing e-mail as spam & Pass ($p=0.945$) \\
    \hline
    7 & E-mail detector passes a meeting-notes e-mail as ham & Pass ($p=0.140$) \\
    \hline
    8 & All three e-mail base learners return a vote & Pass \\
    \hline
    9 & LIME returns word weights for the phishing e-mail & Pass \\
    \hline
    10 & SHAP returns attributions for the phishing e-mail & Pass \\
    \hline
    11 & Application renders without an exception & Pass \\
    \hline
    12 & Clicking \emph{Analyse} renders without an exception & Pass \\
    \hline
    13 & The verdict is shown to the user & Pass \\
    \hline
  \end{tabular}
\end{table}

%------------------------------------------------------------------------------
\chapter{CONCLUSION AND FUTURE WORK}

\section{Conclusion}

Every objective set in Section~\ref{sec:objectives} was met. The hybrid ensemble reached
98.55\% accuracy on SMS and 98.25\% on e-mail against a target of 97\%, and
it matched (to four decimal places) or beat every base learner on ROC-AUC on
both modalities (0.9941
on SMS; an outright best 0.9983 on e-mail). Every single prediction is
delivered with a local LIME explanation over the full ensemble and a SHAP
attribution of the boosted-tree component, and global beeswarm and importance
summaries describe what each model has learned. The complete pipeline ---
corpus download, training, evaluation and figure generation --- reproduces
from a clean checkout in about three minutes on a commodity CPU and is
verified by a 13-check automated test suite, and the four-tab Streamlit
application makes the system operable by a non-expert.

The most instructive result, however, was not a number. The global SHAP
analysis exposed the token \texttt{enron} acting as a shortcut for legitimate
mail --- a genuine dataset bias in the e-mail corpus that deviated from the
naive expectation that a 98\%-accurate model has learned the concept of
``spam''. The bias is a property of the data, not a fault of the pipeline,
but it would have shipped silently inside any black-box deployment. Catching
it required the explanation layer, which is precisely the thesis this
project set out to demonstrate: interpretability is not a presentation
feature bolted onto a classifier, it is a debugging and auditing instrument
that changes what the engineer can see.

The additional validation reported in Chapter~4 strengthens these
conclusions. Paired tests show that the ensemble is significantly better than
both tree models on e-mail and statistically level with Logistic Regression; a
deletion test confirms that LIME's explanations point at the words on which
the decision really depends; the \texttt{enron} shortcut was quantified, with a
single inserted word evading the ensemble on 35.5\% of spam e-mails and
XGBoost alone on 91.7\%; and the heuristic router was identified as the
weakest link of the deployed system. Of 28 realistic messages written for this
project, 26 were classified correctly.

\section{Future work}

\begin{itemize}[itemsep=0pt]
  \item Extend the framework to multilingual and code-mixed corpora, since
        the present datasets are English-only and date from 2002--2012.
  \item Perform cross-domain evaluation (train on one corpus, test on
        another) to quantify and then mitigate the dataset bias exposed by
        SHAP.
  \item Add probability calibration so that the reported spam probabilities
        can be thresholded meaningfully by downstream policy.
  \item Harden the models against adversarially obfuscated spelling (e.g.\
        leet-speak substitutions) and measure robustness explicitly.
  \item Add counterfactual explanations --- ``what is the smallest edit that
        flips this verdict?'' --- as a third, complementary XAI technique.
  \item Benchmark a fine-tuned transformer under the same explanation
        harness to quantify the accuracy--interpretability trade-off
        empirically.
  \item Replace the heuristic router with a learned channel classifier,
        since automatic routing currently lowers pooled accuracy from 98.3\%
        to 92.2\%.
  \item Mask organisation-identity tokens such as \texttt{enron} during
        training and measure how accuracy and robustness change.
  \item Compute the SHAP waterfall in the application directly on the sparse
        feature matrix, as done for the figures in this report, so that the
        interactive explanations match them exactly.
\end{itemize}

%==============================================================================
%  APPENDIX
%==============================================================================
\chapter*{APPENDIX}
\phantomsection\addcontentsline{toc}{chapter}{APPENDIX}

\begin{center}\textbf{\fourteen Appendix-1: Software environment}\end{center}

\vspace{6pt}
\begin{singlespace}
\begin{center}
\begin{tabular}{|c|L{3.4cm}|L{3.6cm}|L{6.0cm}|}
  \hline
  \textbf{S.~No.} & \textbf{Component} & \textbf{Version} & \textbf{Role in the framework} \\
  \hline
  1 & Python & 3.10 or newer (tested on 3.12 and 3.14) & Implementation language of the entire pipeline \\
  \hline
  2 & scikit-learn & 1.4 or newer & Logistic Regression, Random Forest, TF--IDF, metrics, cross-validation \\
  \hline
  3 & XGBoost & 2.0 or newer & Gradient-boosted tree base learner \\
  \hline
  4 & SHAP & 0.45 or newer & Exact TreeSHAP local and global attributions \\
  \hline
  5 & LIME & 0.2.0.1 or newer & Model-agnostic local explanations of the full ensemble \\
  \hline
  6 & pandas / NumPy & 2.2 / 1.26 or newer & Data handling and numerical computation \\
  \hline
  7 & matplotlib & 3.8 or newer & All generated figures (shared publication style) \\
  \hline
  8 & Streamlit & 1.36 or newer & Four-tab interactive demonstration application \\
  \hline
  9 & Git & -- & Version control and collaboration \\
  \hline
\end{tabular}
\end{center}
\end{singlespace}

\vspace{14pt}
\begin{center}\textbf{\fourteen Appendix-2: Engineered stylometric features (complete set)}\end{center}

\vspace{6pt}
\begin{singlespace}
\begin{center}
\begin{tabular}{|c|L{4.6cm}|L{7.4cm}|}
  \hline
  \textbf{S.~No.} & \textbf{Feature} & \textbf{Description} \\
  \hline
  1 & \texttt{CHAR\_COUNT} & Length of the message in characters \\
  \hline
  2 & \texttt{WORD\_COUNT} & Number of words \\
  \hline
  3 & \texttt{AVG\_WORD\_LEN} & Average word length \\
  \hline
  4 & \texttt{DIGIT\_RATIO} & Share of characters that are digits \\
  \hline
  5 & \texttt{UPPERCASE\_RATIO} & Share of letters written in upper case \\
  \hline
  6 & \texttt{ALLCAPS\_WORDS} & Words written fully in capitals (e.g.\ `WINNER') \\
  \hline
  7 & \texttt{PUNCT\_RATIO} & Share of punctuation characters \\
  \hline
  8 & \texttt{NUM\_EXCLAMATIONS} & Count of `!' \\
  \hline
  9 & \texttt{NUM\_QUESTIONS} & Count of `?' \\
  \hline
  10 & \texttt{NUM\_CURRENCY} & Currency symbols and money amounts \\
  \hline
  11 & \texttt{NUM\_URLS} & Links and web addresses detected \\
  \hline
  12 & \texttt{NUM\_EMAILS} & E-mail addresses detected \\
  \hline
  13 & \texttt{NUM\_PHONES} & Phone numbers detected \\
  \hline
  14 & \texttt{SPAM\_KEYWORDS} & Hits in the scam-keyword lexicon (Appendix-4) \\
  \hline
  15 & \texttt{SPECIAL\_CHAR\_RATIO} & Share of non-alphanumeric characters \\
  \hline
\end{tabular}
\end{center}
\end{singlespace}

All fifteen features are computed on the raw message, before normalisation, and
are defined in plain sight in \texttt{src/features.py}, together with the
keyword lexicon in \texttt{src/config.py}, so that they can be audited without
reading model internals.

\vspace{14pt}
\begin{center}\textbf{\fourteen Appendix-3: Derived ensemble weights (normalised CV F1)}\end{center}

\vspace{6pt}
\begin{singlespace}
\begin{center}
\begin{tabular}{|L{3.6cm}|c|c|c|c|}
  \hline
  \textbf{Modality} & \textbf{Base learner} & \textbf{CV F1} & \textbf{Weight $w_m$} & \\
  \cline{1-4}
  \multirow{3}{*}{SMS} & Logistic Regression & 0.9554 & 0.3329 & \multirow{3}{*}{$\sum w_m = 1$} \\
  \cline{2-4}
   & Random Forest & 0.9558 & 0.3331 & \\
  \cline{2-4}
   & XGBoost & 0.9586 & 0.3340 & \\
  \hline
  \multirow{3}{*}{E-mail} & Logistic Regression & 0.9797 & 0.3356 & \multirow{3}{*}{$\sum w_m = 1$} \\
  \cline{2-4}
   & Random Forest & 0.9669 & 0.3311 & \\
  \cline{2-4}
   & XGBoost & 0.9732 & 0.3333 & \\
  \hline
\end{tabular}
\end{center}
\end{singlespace}

The near-uniform weights show that the three learners are comparably strong;
the ensemble's advantage therefore comes from their diversity (linear~/
bagging~/ boosting) rather than from any single dominant model.

\vspace{14pt}
\begin{center}\textbf{\fourteen Appendix-4: Scam-keyword lexicon}\end{center}

\vspace{6pt}
The \texttt{SPAM\_KEYWORDS} feature counts whole-word, case-insensitive matches
of the following 46 words, defined in \texttt{src/config.py}.

\begin{singlespace}
\begin{center}
{\small
\begin{tabular}{|L{2.7cm}|L{2.7cm}|L{2.7cm}|L{2.7cm}|L{2.7cm}|}
  \hline
  free & win & winner & won & winning \\
  \hline
  cash & prize & award & urgent & claim \\
  \hline
  congratulations & congrats & offer & click & subscribe \\
  \hline
  guaranteed & guarantee & credit & loan & lottery \\
  \hline
  million & bonus & discount & cheap & deal \\
  \hline
  limited & expires & expire & alert & verify \\
  \hline
  account & password & bank & refund & investment \\
  \hline
  viagra & pills & ringtone & voucher & gift \\
  \hline
  txt & sms & mobile & selected & reward \\
  \hline
  unsubscribe &  &  &  &  \\
  \hline
\end{tabular}}
\end{center}
\end{singlespace}

\end{document}